\documentclass[11pt,oneside]{article}

% =====================================================================
%  Inverse transseries for the real-argument Fibonacci modulus
%  Research draft for the ProveIt series-and-transseries collection.
% =====================================================================

\usepackage[a4paper,margin=26mm,headheight=15pt]{geometry}
\usepackage[T1]{fontenc}
\usepackage[utf8]{inputenc}
\IfFileExists{libertinus.sty}{\usepackage{libertinus}}{\usepackage{lmodern}}
\usepackage{microtype}
\usepackage{amsmath,amssymb,amsthm,mathtools,mathrsfs,bm}
\usepackage{booktabs,longtable,tabularx,array,multirow}
\usepackage{graphicx}
\usepackage{xcolor}
\usepackage{enumitem}
\usepackage{fancyhdr}
\usepackage{float}
\usepackage{xurl}
\usepackage{listings}
\usepackage[most]{tcolorbox}
\usepackage{hyperref}
\usepackage[nameinlink,noabbrev,capitalise]{cleveref}

\definecolor{linkblue}{RGB}{24,72,124}
\definecolor{deepgreen}{RGB}{23,102,69}
\definecolor{deepred}{RGB}{140,42,42}
\definecolor{deepgold}{RGB}{122,91,19}
\definecolor{shadegray}{RGB}{246,247,249}
\definecolor{rulegray}{RGB}{195,201,208}
\definecolor{palered}{RGB}{251,236,236}
\definecolor{palegold}{RGB}{251,245,230}
\definecolor{palegreen}{RGB}{233,245,243}
\definecolor{paleblue}{RGB}{237,244,251}

\hypersetup{
  hypertexnames=false,
  colorlinks=true,
  linkcolor=linkblue,
  citecolor=deepgreen,
  urlcolor=linkblue,
  pdftitle={Inverting a Real-Argument Fibonacci Function: Log-Periodic Transseries, Exact Coefficients, and a General Product-Reversion Calculus},
  pdfauthor={Research draft prepared for Vladimir Reshetnikov and the ProveIt project},
  pdfsubject={Asymptotic and convergent transseries for the outer inverse branches of an even real interpolation of Fibonacci numbers, with multivariate Lagrange-Good inversion},
  pdfkeywords={Fibonacci numbers, inverse function, transseries, log-periodic asymptotics, complex powers, multivariate Lagrange inversion, Good inversion, generalized binomial coefficients}
}

\pagestyle{fancy}
\fancyhf{}
\fancyhead[L]{\small Inverse Fibonacci transseries}
\fancyhead[R]{\small\nouppercase{\leftmark}}
\fancyfoot[C]{\thepage}
\renewcommand{\headrulewidth}{0.4pt}

\setcounter{tocdepth}{2}
\setcounter{secnumdepth}{3}
\setlength{\parindent}{0pt}
\setlength{\parskip}{0.48em}
\setlist[itemize]{leftmargin=1.45em,itemsep=0.25ex,topsep=0.55ex}
\setlist[enumerate]{leftmargin=1.85em,itemsep=0.35ex,topsep=0.55ex}
\allowdisplaybreaks[2]
\numberwithin{equation}{section}

\theoremstyle{plain}
\newtheorem{theorem}{Theorem}[section]
\newtheorem{proposition}[theorem]{Proposition}
\newtheorem{lemma}[theorem]{Lemma}
\newtheorem{corollary}[theorem]{Corollary}
\theoremstyle{definition}
\newtheorem{definition}[theorem]{Definition}
\newtheorem{example}[theorem]{Example}
\newtheorem{algorithm}[theorem]{Algorithm}
\theoremstyle{remark}
\newtheorem{remark}[theorem]{Remark}
\newtheorem{warning}[theorem]{Warning}

\newtcolorbox{mainresultbox}[1][]{%
  enhanced,breakable,colback=paleblue,colframe=linkblue,
  boxrule=0.75pt,arc=1.2mm,left=2.2mm,right=2.2mm,top=1.5mm,bottom=1.5mm,
  title={Main result},fonttitle=\bfseries,#1}
\newtcolorbox{summarybox}[1][]{%
  enhanced,breakable,colback=palegold,colframe=deepgold,
  boxrule=0.6pt,arc=1.2mm,left=1.8mm,right=1.8mm,top=1.2mm,bottom=1.2mm,
  title={Section summary},fonttitle=\bfseries,#1}
\newtcolorbox{checkpointbox}[1][]{%
  enhanced,breakable,colback=palegreen,colframe=deepgreen,
  boxrule=0.6pt,arc=1.2mm,left=1.8mm,right=1.8mm,top=1.2mm,bottom=1.2mm,
  title={Structural checkpoint},fonttitle=\bfseries,#1}
\newtcolorbox{warningbox}[1][]{%
  enhanced,breakable,colback=palered,colframe=deepred,
  boxrule=0.6pt,arc=1.2mm,left=1.8mm,right=1.8mm,top=1.2mm,bottom=1.2mm,
  title={A point requiring care},fonttitle=\bfseries,#1}

\lstdefinestyle{wlstyle}{%
  language=Mathematica,
  basicstyle=\ttfamily\small,columns=fullflexible,frame=single,
  rulecolor=\color{rulegray},backgroundcolor=\color{shadegray},
  showstringspaces=false,keepspaces=true,breaklines=true,
  xleftmargin=0.4em,xrightmargin=0.4em,aboveskip=0.8em,belowskip=0.8em}
\lstset{style=wlstyle}

\newcommand{\Fcal}{\mathcal{F}}
\newcommand{\e}{\mathrm e}
\newcommand{\ii}{\mathrm i}
\newcommand{\dd}{\,\mathrm d}
\newcommand{\Ord}{\mathcal O}
\newcommand{\littleo}{\mathrm o}
\newcommand{\phig}{\varphi}
\newcommand{\Logp}{\operatorname{Log}}
\newcommand{\Res}{\operatorname*{Res}}
\newcommand{\Coeff}[2]{[\,#1^{#2}\,]}
\newcommand{\multi}[1]{\bm{#1}}
\newcommand{\Nzero}{\mathbb N_0}
\newcommand{\Talg}{\mathscr A_{\rho}}
\newcommand{\C}{\mathbb C}
\newcommand{\R}{\mathbb R}
\newcommand{\mean}[1]{\left\langle #1\right\rangle}
\DeclareMathOperator{\arccsch}{arccsch}
\DeclareMathOperator{\arsinh}{arsinh}
\DeclareMathOperator{\val}{val}
\DeclareMathOperator{\supp}{supp}

\title{\textbf{Inverting a Real-Argument Fibonacci Function}\\[0.35em]
\Large Log-Periodic Transseries, Exact Multivariate Coefficients,\\
and a General Product-Reversion Calculus}
\author{Research draft prepared for Vladimir Reshetnikov\\and the ProveIt project}
\date{3 September 2026}

\begin{document}
\maketitle

\begin{abstract}
We study the nonnegative even interpolation
\[
 \Fcal(x)=\sqrt{\frac25}\,
 \sqrt{\cosh\!\bigl(2x\,\arccsch 2\bigr)-\cos(\pi x)},
 \qquad x\in\R,
\]
which satisfies \(\Fcal(n)=F_n\) for every nonnegative integer \(n\).
The map is not globally one-to-one, but it is strictly increasing on
\([2,\infty)\).  Hence it has a positive outer inverse \(X_+(y)\) for
\(y\ge1\), and evenness gives the negative outer inverse
\(X_-(y)=-X_+(y)\).  We determine their complete large-\(y\) transseries.

Put \(\lambda=2\log\varphi\), \(\rho=\pi/\lambda\),
\(R=5y^2\), and \(u=\e^{-\lambda x}\).  The inverse equation admits the
factorization
\[
 Ru=(1-u^{1-\ii\rho})(1-u^{1+\ii\rho}).
\]
This identity converts the inverse problem into an analytic bivariate
reversion at the ordinary origin.  Multivariate Lagrange--Good inversion
then gives the explicit all-orders formula
\[
 X_+(y)=\frac{\log R}{\lambda}
 +\sum_{\substack{j,k\ge0\\j+k\ge1}}
 \frac{(-1)^{j+k+1}}{\lambda s_{jk}}
 \binom{s_{jk}}j\binom{s_{jk}}k R^{-s_{jk}},
 \qquad
 s_{jk}=j+k-\ii\rho(j-k).
\]
The series is not merely formal: it converges for all sufficiently large
\(y\).  Pairing conjugate terms reorganizes it as
\[
 X_+(y)=\xi+\sum_{n\ge1}D_n(\theta)R^{-n},
 \qquad
 \xi=\frac{\log R}{\lambda},\quad \theta=\pi\xi,
\]
where each \(D_n\) is an explicit real trigonometric polynomial whose
harmonics have the same parity as \(n\).  We derive the first sectors,
a triangular recurrence, uniform remainder bounds, differentiation rules,
a hypergeometric--elliptic formula for the angular mean, and the exact
Banach-space radius of uniform convergence.  Numerical tests include the
integer checkpoints \(X_+(2)=3\) and \(X_+(5)=5\).

The second part develops a reusable theory for inverses of finite products
\[
 R=\e^{\lambda x}\prod_{r=1}^m
 (1-c_r\e^{-\lambda\alpha_r x})^{\beta_r},
 \qquad \Re\alpha_r>0,
\]
obtaining a closed multi-index coefficient formula.  We then extend the
method to arbitrary analytic perturbations of a dominant exponential and
to mixed carriers such as \(\e^{\lambda x}x^\mu\), where Lambert \(W\)
and polynomial--logarithmic coefficient fields enter.  The resulting
framework explains systematically why conjugate complex exponents generate
log-periodic or quasiperiodic inverse transseries.
\end{abstract}

\begin{mainresultbox}
Let
\[
 \lambda=2\log\varphi,\qquad \rho=\frac{\pi}{\lambda},\qquad
 R=5y^2,
\]
and define
\[
 \xi=\frac{\log R}{\lambda}=\log_{\varphi}(\sqrt5\,y),
 \qquad \theta=\rho\log R=\pi\xi,
 \qquad q=R^{-1}.
\]
For the positive outer branch, as \(y\to+\infty\),
\begin{align*}
 X_+(y)={}&\xi+\frac{2q}{\lambda}\cos\theta\\
 &-\frac{q^2}{\lambda}
   \bigl(2+\cos2\theta+2\rho\sin2\theta\bigr)\\
 &+\frac{q^3}{\lambda}
 \left[
   (6-\rho^2)\cos\theta+5\rho\sin\theta
   +\left(\frac23-3\rho^2\right)\cos3\theta
   +3\rho\sin3\theta
 \right]
 +\Ord(q^4).
\end{align*}
Equivalently, the leading correction is
\[
 X_+(y)=\log_{\varphi}(\sqrt5\,y)
 +\frac{\cos\!\bigl(\pi\log_{\varphi}(\sqrt5\,y)\bigr)}
 {5(\log\varphi)y^2}
 +\Ord(y^{-4}).
\]
The complete coefficient formula is
\[
 X_+(y)=\frac{\log R}{\lambda}
 +\sum_{\substack{j,k\ge0\\j+k\ge1}}
 \frac{(-1)^{j+k+1}}{\lambda s_{jk}}
 \binom{s_{jk}}j\binom{s_{jk}}k R^{-s_{jk}},
 \qquad s_{jk}=j+k-\ii\rho(j-k).
\]
For every fixed \(N\), truncation at \(j+k=N\) has error
\(\Ord(R^{-N-1})=\Ord(y^{-2N-2})\).  The negative outer branch is
\(X_-(y)=-X_+(y)\).
\end{mainresultbox}

\tableofcontents
\newpage
\section{Introduction}

The ordinary Fibonacci sequence is discrete, while inversion is naturally a
continuous or complex-analytic operation.  There are many inequivalent ways
to continue Fibonacci numbers away from the integers.  The continuation in
this article is distinguished by three elementary features: it is real and
nonnegative on the real axis, it is even, and it agrees exactly with the
nonnegative Fibonacci numbers at nonnegative integers.  These features are
purchased at the cost of global injectivity and, at the origin, ordinary
analyticity.  Neither issue obstructs inversion at large values.

The asymptotic problem has an unusual scale.  The dominant exponential gives
\(
 x\sim\log_{\varphi}(\sqrt5\,y)
\), but the first correction is not a power of \(1/\log y\).  It has size
\(y^{-2}\) and carries a periodic function of \(\log y\).  Thus it is beyond
all inverse-logarithmic orders while still algebraic in \(y^{-1}\).  In the
language of transseries, the relevant monomials are complex powers
\[
 R^{-n+\ii\rho h}=R^{-n}\e^{\ii h\rho\log R},
 \qquad n\in\mathbb N,
\]
organized into finite Fourier packets at every real order \(R^{-n}\).
Conjugation makes the final answer real.

A direct Taylor reversion of the original square-root formula is possible,
but it hides the decisive structure.  Squaring and introducing
\(u=\e^{-\lambda x}\) reveals an exact factorization into two conjugate
factors.  The inverse then becomes a bivariate fixed-point problem with an
ordinary analytic germ at the origin.  This makes the multivariate
Lagrange--Good theorem almost perfectly adapted to the problem and produces
a closed coefficient formula.

The article has two goals.  First, it gives a complete treatment of the
specific Fibonacci interpolation: branches, explicit sectors, convergence,
error certification, computation, and numerical checks.  Second, it isolates
the mechanism and proves a multi-factor inversion theorem applicable to a
large class of exponential products.  Arbitrary analytic perturbations and
mixed dominant carriers are treated afterward.

\section{The real-argument Fibonacci modulus}

\subsection{Constants and equivalent forms}

Let
\begin{equation}\label{eq:phi-lambda-rho}
 \varphi=\frac{1+\sqrt5}{2},\qquad
 a=\arccsch 2,\qquad
 \lambda=2a,
 \qquad \omega=\pi,
 \qquad \rho=\frac{\omega}{\lambda}.
\end{equation}
Because \(\sinh a=1/2\), the positive number \(t=\e^a\) satisfies
\(t-t^{-1}=1\), hence \(t=\varphi\).  Therefore
\begin{equation}\label{eq:arccsch-phi}
 a=\log\varphi,
 \qquad
 \lambda=2\log\varphi.
\end{equation}
Numerically,
\[
 \lambda=0.9624236501192068949955\ldots,
 \qquad
 \rho=3.2642513026364969\ldots.
\]

Define
\begin{equation}\label{eq:F-definition}
 \Fcal(x)=\sqrt{\frac25}
 \sqrt{\cosh(\lambda x)-\cos(\omega x)},
 \qquad x\in\R,
\end{equation}
where the nonnegative square root is understood.  Squaring and expanding the
hyperbolic cosine gives
\begin{equation}\label{eq:F-square-exp}
 \Fcal(x)^2
 =\frac15\left(\e^{\lambda x}+\e^{-\lambda x}
                   -2\cos(\omega x)\right).
\end{equation}
Since \(\lambda=2\log\varphi\), this can also be written as the modulus of a
complex Binet interpolation:
\begin{equation}\label{eq:modulus-binet}
 \Fcal(x)
 =\frac1{\sqrt5}\left|
    \varphi^x-\e^{\ii\pi x}\varphi^{-x}
  \right|.
\end{equation}
The modulus in \cref{eq:modulus-binet} is essential: the complex expression
inside it is an analytic continuation, whereas \(\Fcal\) itself is an even,
nonnegative real function.

\begin{proposition}[Interpolation of Fibonacci numbers]\label{prop:integer-interpolation}
For every integer \(n\ge0\),
\[
 \Fcal(n)=F_n,
\]
where \(F_0=0\), \(F_1=1\), and \(F_{n+2}=F_{n+1}+F_n\).
\end{proposition}

\begin{proof}
Binet's formula is
\[
 F_n=\frac{\varphi^n-(-\varphi^{-1})^n}{\sqrt5}.
\]
It is nonnegative for \(n\ge0\), and therefore
\begin{align*}
 F_n^2
 &=\frac15\left(\varphi^{2n}+\varphi^{-2n}-2(-1)^n\right)\\
 &=\frac15\left(\e^{\lambda n}+\e^{-\lambda n}
                  -2\cos(\pi n)\right)
 =\Fcal(n)^2.
\end{align*}
Taking nonnegative square roots proves the claim.
\end{proof}

\begin{remark}
The extension is not the usual signed extension to negative indices.  It is
even, so \(\Fcal(-n)=F_n\), whereas the discrete signed convention gives
\(F_{-n}=(-1)^{n+1}F_n\).  This distinction matters when the two outer
inverse branches are named.
\end{remark}

\subsection{Basic geometry and the outer inverse branches}

Equation \cref{eq:F-definition} immediately shows that \(\Fcal\) is even and
that its only real zero is \(x=0\).  Near the origin,
\begin{equation}\label{eq:F-origin}
 \Fcal(x)=\sqrt{\frac{\lambda^2+\omega^2}{5}}\,|x|
 +\Ord(|x|^3),
\end{equation}
so the even interpolation has a cusp there.

For \(x>0\), differentiation of \cref{eq:F-square-exp} gives
\begin{equation}\label{eq:F-derivative}
 \Fcal'(x)=
 \frac{\lambda\sinh(\lambda x)+\omega\sin(\omega x)}
      {5\Fcal(x)}.
\end{equation}
The oscillatory sine can create a small nonmonotone region, but exponential
growth eventually dominates it.

\begin{proposition}[Monotonicity of the outer branch]\label{prop:outer-monotone}
Let
\begin{equation}\label{eq:x-star}
 x_*=\frac1\lambda\arsinh\!\left(\frac\omega\lambda\right)
 =1.972992734136994\ldots.
\end{equation}
Then \(\Fcal'(x)>0\) for every \(x>x_*\).  In particular,
\(\Fcal\) is strictly increasing on \([2,\infty)\), and it maps that
interval bijectively onto \([1,\infty)\).
\end{proposition}

\begin{proof}
By \cref{eq:F-derivative}, the sign is the sign of its numerator.  Since
\(\sin(\omega x)\ge-1\),
\[
 \lambda\sinh(\lambda x)+\omega\sin(\omega x)
 \ge \lambda\sinh(\lambda x)-\omega.
\]
The right-hand side is positive exactly when \(x>x_*\).  Moreover,
\(x_*<2\) and \(\Fcal(2)=F_2=1\).  Since \(\Fcal(x)\to\infty\), the stated
bijection follows.
\end{proof}

\begin{definition}[Outer inverse branches]\label{def:outer-branches}
For \(y\ge1\), let \(X_+(y)\in[2,\infty)\) be the unique solution of
\(\Fcal(x)=y\).  Define
\[
 X_-(y)=-X_+(y).
\]
These are the positive and negative outer inverse branches.
\end{definition}

On the entire real line, \(\Fcal\) is not globally one-to-one.  Numerically,
the two nontrivial critical points in \((1,2)\) are approximately
\[
 1.132379068\ldots,
 \qquad
 1.719250886\ldots,
\]
with function values \(1.013859885\ldots\) and
\(0.911172684\ldots\), respectively.  Consequently, some moderate values
have additional inner preimages.  This finite-scale geometry disappears
from the large-value problem: because \([0,2]\) is compact, every
sufficiently large \(y\) has exactly the two real preimages \(X_+(y)\) and
\(X_-(y)\).

\begin{warningbox}
The phrase ``the inverse of \(\Fcal\)'' is ambiguous unless a branch is
specified.  Every asymptotic formula below first refers to \(X_+\); the
negative real branch follows by evenness.  Complex branches require an
additional logarithm index and are discussed in \cref{sec:complex-branches}.
\end{warningbox}

\section{Exact normalization and factorization}

\subsection{The squared inverse equation}

Fix the positive outer branch and put
\begin{equation}\label{eq:R-def}
 R=5y^2.
\end{equation}
The equation \(\Fcal(x)=y\) becomes
\begin{equation}\label{eq:R-equation-x}
 R=\e^{\lambda x}+\e^{-\lambda x}-2\cos(\omega x).
\end{equation}
For positive real \(x\), introduce
\begin{equation}\label{eq:u-def}
 u=\e^{-\lambda x}\in(0,1),
 \qquad \rho=\frac\omega\lambda.
\end{equation}
Using the real logarithm of \(u\),
\[
 \e^{\ii\omega x}=u^{-\ii\rho},
 \qquad
 \e^{-\ii\omega x}=u^{\ii\rho}.
\]
Therefore \cref{eq:R-equation-x} is equivalent to
\begin{equation}\label{eq:R-u-unfactored}
 R=u^{-1}+u-u^{\ii\rho}-u^{-\ii\rho}.
\end{equation}
Multiplication by \(u\) produces the key identity.

\begin{theorem}[Conjugate factorization]\label{thm:factorization}
With
\begin{equation}\label{eq:alpha-def}
 \alpha=1-\ii\rho,
 \qquad
 \bar\alpha=1+\ii\rho,
\end{equation}
the positive inverse equation is exactly
\begin{equation}\label{eq:factorized-u}
 Ru=(1-u^\alpha)(1-u^{\bar\alpha}).
\end{equation}
\end{theorem}

\begin{proof}
Expanding the right-hand side gives
\[
 (1-u^{1-\ii\rho})(1-u^{1+\ii\rho})
 =1-u^{1-\ii\rho}-u^{1+\ii\rho}+u^2,
\]
which is precisely \(u\) times the right-hand side of
\cref{eq:R-u-unfactored}.
\end{proof}

This factorization is the structural reason the inverse coefficients admit
a closed formula.  Without it one still obtains a recursion, but the two
oscillatory pieces appear unrelated.  In factorized form they are a
conjugate pair.

\subsection{Baseline, phase, and the normalized implicit equation}

Define
\begin{equation}\label{eq:q-xi-theta}
 q=R^{-1},
 \qquad
 \xi=\frac{\log R}{\lambda},
 \qquad
 \theta=\rho\log R=\omega\xi.
\end{equation}
The dominant balance in \cref{eq:R-equation-x} is
\(R\sim\e^{\lambda x}\), hence \(x\sim\xi\).  Write
\begin{equation}\label{eq:delta-def}
 x=\xi+\delta.
\end{equation}
Since \(\e^{\lambda\xi}=R\), division of
\cref{eq:R-equation-x} by \(R\) yields the exact scalar equation
\begin{equation}\label{eq:H-equation}
 H(\delta;q,\theta):=
 \e^{\lambda\delta}
 +q^2\e^{-\lambda\delta}
 -2q\cos(\theta+\omega\delta)-1=0.
\end{equation}
At \(q=0\), this reduces to \(\e^{\lambda\delta}=1\), and the real germ is
\(\delta=0\).  Furthermore,
\[
 \frac{\partial H}{\partial\delta}(0;0,\theta)=\lambda\ne0.
\]
The analytic implicit-function theorem therefore gives a unique small
solution for \(|q|\) sufficiently small, uniformly in real \(\theta\).
Equation \cref{eq:H-equation} already proves the existence of an expansion
in powers of \(q\) whose coefficients are trigonometric polynomials.  The
factorization will do much more: it will calculate every coefficient in one
line.

\subsection{Bivariate lift}

From \cref{eq:factorized-u}, set
\begin{equation}\label{eq:VABXY}
 V=\frac uq,
 \qquad
 A=u^\alpha,
 \qquad
 B=u^{\bar\alpha},
 \qquad
 X=q^\alpha,
 \qquad
 Y=q^{\bar\alpha}.
\end{equation}
Then
\begin{equation}\label{eq:V-product}
 V=(1-A)(1-B),
\end{equation}
and, because \(u=qV\),
\begin{equation}\label{eq:AB-fixed-point}
 A=XV^\alpha,
 \qquad
 B=YV^{\bar\alpha}.
\end{equation}
Finally,
\begin{equation}\label{eq:delta-logV}
 \delta=-\frac1\lambda\log V.
\end{equation}
Equations \cref{eq:V-product,eq:AB-fixed-point,eq:delta-logV} are an
ordinary bivariate analytic reversion problem at \((X,Y)=(0,0)\).
For positive real \(q\),
\begin{equation}\label{eq:X-Y-phase}
 X=q\e^{\ii\theta},
 \qquad
 Y=q\e^{-\ii\theta}.
\end{equation}
Thus every monomial \(X^jY^k\) becomes
\(q^{j+k}\e^{\ii(j-k)\theta}\).

\begin{checkpointbox}
The large variable \(y\) has disappeared from the analytic reversion itself.
It enters only through the small radial variable \(q=(5y^2)^{-1}\) and the
logarithmic angle \(\theta=\rho\log(5y^2)\).  The inverse correction is an
analytic bivariate germ restricted to the spiral
\((X,Y)=(q\e^{\ii\theta},q\e^{-\ii\theta})\).
\end{checkpointbox}

\section{The log-periodic transseries algebra}\label{sec:algebra}

\subsection{Complex powers and Fourier packets}

For the present problem the natural monomials are
\begin{equation}\label{eq:monomial-R}
 R^{-s_{jk}},
 \qquad
 s_{jk}=\alpha j+\bar\alpha k
       =j+k-\ii\rho(j-k),
 \qquad j,k\in\Nzero.
\end{equation}
Their absolute sizes are \(R^{-(j+k)}\).  If
\begin{equation}\label{eq:n-h}
 n=j+k,
 \qquad h=j-k,
\end{equation}
then \(|h|\le n\), \(h\equiv n\pmod2\), and
\begin{equation}\label{eq:monomial-qtheta}
 R^{-s_{jk}}=q^n\e^{\ii h\theta}.
\end{equation}
At each fixed radial order \(q^n\), only finitely many Fourier harmonics are
possible.

\begin{definition}[Fibonacci log-periodic algebra]\label{def:Talg}
Let \(\Talg\) be the set of formal series
\begin{equation}\label{eq:Talg-form}
 T(q,\theta)=\sum_{n\ge0}D_n(\theta)q^n,
\end{equation}
where
\begin{equation}\label{eq:Dn-cone}
 D_n(\theta)=
 \sum_{\substack{|h|\le n\\h\equiv n\ (2)}}d_{n,h}\e^{\ii h\theta}
\,
\end{equation}
and, for real series, \(d_{n,-h}=\overline{d_{n,h}}\).  The valuation of a
nonzero series is the least \(n\) for which \(D_n\ne0\).
\end{definition}

The cone and parity conditions in \cref{eq:Dn-cone} are closed under
multiplication.  The resulting \(q\)-adic metric makes \(\Talg\) complete.
If \(T\) has positive valuation, then analytic substitutions such as
\(\exp T\), \(\log(1+T)\), and \(\cos(\theta+T)\) are well defined by their
ordinary Taylor series: only finitely many terms contribute to any prescribed
radial order.  Consequently \cref{eq:H-equation} has a unique formal solution
\(\delta\in q\Talg\), independently of convergence questions.

\subsection{Differentiation}

The algebra is also stable under the Euler derivative.  Since
\(q=R^{-1}\) and \(\theta=\rho\log R\),
\begin{equation}\label{eq:euler-action}
 R\frac{\dd}{\dd R}
 \left(q^n\e^{\ii h\theta}\right)
 =(-n+\ii\rho h)q^n\e^{\ii h\theta}
 =-s_{jk}R^{-s_{jk}}.
\end{equation}
In the original variable \(y\), because \(R=5y^2\),
\begin{equation}\label{eq:y-euler-action}
 y\frac{\dd}{\dd y}
 \left(q^n\e^{\ii h\theta}\right)
 =2(-n+\ii\rho h)q^n\e^{\ii h\theta}.
\end{equation}
Thus differentiation does not change radial valuation.  Ordinary derivatives
add the expected powers of \(y^{-1}\).

\subsection{Why this is a transseries rather than an ordinary power series}

If \(q\) and \(\theta\) were independent, \cref{eq:Talg-form} would simply
be a power series with trigonometric-polynomial coefficients.  Along the
physical curve, however,
\(
 \theta=-\rho\log q
\), so every Fourier factor is a complex power of \(q\).
The scale is therefore a finitely generated grid of complex powers.  It is
well based by the real valuation \(\Re s_{jk}=j+k\), even though the complex
exponents are not linearly ordered by their imaginary parts.

Relative to the coarse expansion scale generated by
\(\log y\), \(\log\log y\), and inverse logarithms, every nonconstant sector
of \(\Talg\) is beyond all orders:
\[
 y^{-2}=\littleo\!\left((\log y)^{-M}\right)
 \qquad\text{for every fixed }M>0.
\]
This explains why a purely polynomial--logarithmic analysis recovers only
the carrier \(\xi\) and misses the oscillatory corrections.

\begin{summarybox}
The correct coefficient ring is not merely \(\R[[y^{-2}]]\).  It is a
complete valued algebra of Fourier packets in \(\log y\).  The bivariate
variables \(X\) and \(Y\) linearize its support: \(X^jY^k\) corresponds to
the complex-power monomial \(R^{-s_{jk}}\).
\end{summarybox}
\section{Multivariate Lagrange--Good inversion}\label{sec:good}

\subsection{The coefficient theorem}

We use the following form of multivariate Lagrange inversion.  For a
multi-index \(\multi n=(n_1,\ldots,n_m)\), write
\(
 \multi Q^{\multi n}=\prod_r Q_r^{n_r}
\) and similarly for \(\multi A^{\multi n}\).

\begin{theorem}[Lagrange--Good coefficient formula]\label{thm:good}
Suppose analytic germs \(A_r=A_r(\multi Q)\) are determined near the origin
by
\begin{equation}\label{eq:Good-system}
 A_r=Q_r\Phi_r(\multi A),
 \qquad r=1,\ldots,m,
 \qquad \Phi_r(\multi0)\ne0.
\end{equation}
Then for every analytic germ \(f\),
\begin{equation}\label{eq:Good-formula}
 [\multi Q^{\multi n}]f(\multi A(\multi Q))
 = [\multi A^{\multi n}]
 f(\multi A)\prod_{r=1}^m\Phi_r(\multi A)^{n_r}
 \Delta(\multi A),
\end{equation}
where
\begin{equation}\label{eq:Good-determinant}
 \Delta(\multi A)=
 \det\!\left(
   \delta_{rs}-A_s\frac{\partial}{\partial A_s}
                   \log\Phi_r(\multi A)
 \right)_{r,s=1}^m.
\end{equation}
\end{theorem}

The theorem can be proved by a multivariate residue change of variables.
The determinant is the Jacobian correction for
\(Q_r=A_r/\Phi_r(\multi A)\).  We will not merely invoke the formula: in the
present system its determinant has rank one, and the entire coefficient
calculation collapses to generalized binomial coefficients.

\subsection{Evaluation for the Fibonacci system}

For \cref{eq:AB-fixed-point},
\[
 \Phi_1=V^\alpha,
 \qquad
 \Phi_2=V^{\bar\alpha},
 \qquad
 V=(1-A)(1-B).
\]
Because
\[
 \frac{\partial\log V}{\partial A}=-\frac1{1-A},
 \qquad
 \frac{\partial\log V}{\partial B}=-\frac1{1-B},
\]
the Good matrix is
\[
 \begin{pmatrix}
  1+\dfrac{\alpha A}{1-A} & \dfrac{\alpha B}{1-B}\\[0.8em]
  \dfrac{\bar\alpha A}{1-A} & 1+\dfrac{\bar\alpha B}{1-B}
 \end{pmatrix}.
\]
It is an identity matrix plus a rank-one matrix, so
\begin{equation}\label{eq:Delta-fib}
 \Delta(A,B)=1+\frac{\alpha A}{1-A}
                +\frac{\bar\alpha B}{1-B}.
\end{equation}

For \(j,k\ge0\), define
\begin{equation}\label{eq:sjk-repeat}
 s_{jk}=\alpha j+\bar\alpha k.
\end{equation}
Introduce the auxiliary coefficient
\begin{equation}\label{eq:Pjk-def}
 P_{jk}(t)=[A^jB^k]V^t\Delta(A,B).
\end{equation}
The next lemma is the computational core.

\begin{lemma}[Binomial collapse]\label{lem:binomial-collapse}
For \((j,k)\ne(0,0)\),
\begin{equation}\label{eq:Pjk-formula}
 P_{jk}(t)=(-1)^{j+k}
 \binom tj\binom tk\left(1-\frac{s_{jk}}t\right),
\end{equation}
where
\begin{equation}\label{eq:generalized-binomial}
 \binom zn=\frac{z(z-1)\cdots(z-n+1)}{n!},
 \qquad \binom z0=1.
\end{equation}
The apparent singularity at \(t=0\) in \cref{eq:Pjk-formula} is removable
when the expression is interpreted as the polynomial obtained from
\cref{eq:Pjk-def}.
\end{lemma}

\begin{proof}
The contribution of the constant term \(1\) in \cref{eq:Delta-fib} is
\[
 [A^jB^k](1-A)^t(1-B)^t
 =(-1)^{j+k}\binom tj\binom tk.
\]
For the \(A\)-term in the determinant, the coefficient is
\begin{align*}
 &\alpha[A^{j-1}](1-A)^{t-1}[B^k](1-B)^t\\
 &\quad=\alpha(-1)^{j+k-1}
          \binom{t-1}{j-1}\binom tk
 =-(-1)^{j+k}\binom tj\binom tk\frac{\alpha j}{t}.
\end{align*}
The \(B\)-term similarly contributes the factor
\(-\bar\alpha k/t\).  Adding the three pieces gives
\cref{eq:Pjk-formula}.
\end{proof}

\subsection{Differentiating at the forced zero}

Apply \cref{thm:good} to the auxiliary function \(f(A,B)=V^\tau\).  Since
\(\Phi_1^j\Phi_2^k=V^{s_{jk}}\),
\begin{equation}\label{eq:Vtau-coeff}
 [X^jY^k]V(A(X,Y),B(X,Y))^\tau
 =P_{jk}(\tau+s_{jk}).
\end{equation}
Differentiating at \(\tau=0\) gives
\begin{equation}\label{eq:logV-coeff-P}
 [X^jY^k]\log V=P_{jk}'(s_{jk}).
\end{equation}
The factor \(1-s_{jk}/t\) in \cref{eq:Pjk-formula} vanishes at
\(t=s_{jk}\), so derivatives of the binomial factors disappear and
\begin{equation}\label{eq:Pprime}
 P_{jk}'(s_{jk})=
 \frac{(-1)^{j+k}}{s_{jk}}
 \binom{s_{jk}}j\binom{s_{jk}}k.
\end{equation}
Combining this with \cref{eq:delta-logV} proves the explicit coefficient
formula.

\begin{theorem}[Exact bivariate inverse germ]\label{thm:bivariate-inverse}
In a neighborhood of \((X,Y)=(0,0)\), the inverse correction is the
convergent bivariate series
\begin{equation}\label{eq:delta-XY}
 \delta(X,Y)=
 \sum_{\substack{j,k\ge0\\j+k\ge1}}
 C_{jk}X^jY^k,
\end{equation}
where
\begin{equation}\label{eq:Cjk}
 C_{jk}=\frac{(-1)^{j+k+1}}{\lambda s_{jk}}
 \binom{s_{jk}}j\binom{s_{jk}}k,
 \qquad
 s_{jk}=j+k-\ii\rho(j-k).
\end{equation}
\end{theorem}

\begin{proof}
The analytic implicit-function theorem applied to
\cref{eq:AB-fixed-point} gives analytic germs \(A(X,Y)\) and \(B(X,Y)\),
so \(V\) and \(\delta\) are analytic.  Equations
\cref{eq:logV-coeff-P,eq:Pprime,eq:delta-logV} calculate every nonconstant
Taylor coefficient.  The constant term vanishes because \(V(0,0)=1\).
\end{proof}

\begin{remark}[Why the denominator is harmless]
For \((j,k)\ne(0,0)\),
\(
 \Re s_{jk}=j+k>0
\), hence \(s_{jk}\ne0\).  No resonance or small-denominator problem occurs
in this positive cone.  The imaginary part records phase; it cannot cancel
the positive real valuation.
\end{remark}

\section{Complete inverse transseries}\label{sec:complete-transseries}

\subsection{Complex-power form}

Substitute \cref{eq:X-Y-phase} into \cref{eq:delta-XY}.  Since
\(X^jY^k=R^{-s_{jk}}\), we obtain the central result.

\begin{theorem}[All-orders inverse of the Fibonacci modulus]\label{thm:all-orders}
Let \(X_+\) be the outer inverse from \cref{def:outer-branches}.  For all
sufficiently large positive \(y\),
\begin{equation}\label{eq:all-orders-complex}
 \boxed{
 X_+(y)=\frac{\log R}{\lambda}
 +\sum_{\substack{j,k\ge0\\j+k\ge1}}
 \frac{(-1)^{j+k+1}}{\lambda s_{jk}}
 \binom{s_{jk}}j\binom{s_{jk}}k R^{-s_{jk}}
 }
\end{equation}
with
\begin{equation}\label{eq:R-lambda-rho-again}
 R=5y^2,
 \qquad \lambda=2\log\varphi,
 \qquad \rho=\frac\pi\lambda,
 \qquad s_{jk}=j+k-\ii\rho(j-k).
\end{equation}
The series converges to the exact branch.  Independently of convergence,
it is a complete Poincare asymptotic expansion: for every fixed \(N\),
\begin{equation}\label{eq:truncation-complex}
 X_+(y)=\frac{\log R}{\lambda}
 +\sum_{1\le j+k\le N}C_{jk}R^{-s_{jk}}
 +\Ord(R^{-N-1}).
\end{equation}
The negative branch satisfies
\begin{equation}\label{eq:negative-branch}
 X_-(y)=-X_+(y).
\end{equation}
\end{theorem}

\begin{proof}
The bivariate series in \cref{thm:bivariate-inverse} converges in a
polydisc about the origin.  Along the physical curve
\((X,Y)=(q\e^{\ii\theta},q\e^{-\ii\theta})\), both variables tend to zero
as \(y\to\infty\), so the series converges and solves the exact normalized
equation.  Uniqueness of the small real solution identifies it with the
outer branch.  Truncation by total degree leaves an analytic remainder of
total degree at least \(N+1\); because \(|X|=|Y|=q=R^{-1}\), this is
\(\Ord(R^{-N-1})\).  Evenness proves \cref{eq:negative-branch}.
\end{proof}

\subsection{Real Fourier-sector form}

The coefficients satisfy
\begin{equation}\label{eq:C-conjugation}
 C_{kj}=\overline{C_{jk}}.
\end{equation}
Consequently terms with \((j,k)\) and \((k,j)\) pair to a real expression.
For \(n\ge1\) and admissible \(h\), define
\begin{equation}\label{eq:dhat-nh}
 \widehat D_{n,h}=
 \frac{(-1)^{n+1}}{\lambda(n-\ii\rho h)}
 \binom{n-\ii\rho h}{(n+h)/2}
 \binom{n-\ii\rho h}{(n-h)/2}.
\end{equation}
Then put
\begin{equation}\label{eq:Dn-definition}
 D_n(\theta)=
 \sum_{\substack{|h|\le n\\h\equiv n\ (2)}}
 \widehat D_{n,h}\e^{\ii h\theta}.
\end{equation}

\begin{corollary}[Real log-periodic sectors]\label{cor:real-sectors}
The positive outer inverse can be written
\begin{equation}\label{eq:all-orders-real}
 \boxed{
 X_+(y)=\xi+\sum_{n\ge1}D_n(\theta)q^n,
 \qquad
 q=(5y^2)^{-1},\quad
 \xi=\log_{\varphi}(\sqrt5\,y),\quad
 \theta=\pi\xi.
 }
\end{equation}
Each \(D_n\) is real for real \(\theta\), has Fourier degree at most \(n\),
and contains only harmonics congruent to \(n\) modulo \(2\).  Moreover,
\begin{equation}\label{eq:D-parity}
 D_n(\theta+\pi)=(-1)^nD_n(\theta).
\end{equation}
\end{corollary}

\begin{proof}
Set \(n=j+k\) and \(h=j-k\) in \cref{eq:all-orders-complex}.
Equation \cref{eq:C-conjugation} implies
\(
 \widehat D_{n,-h}=\overline{\widehat D_{n,h}}
\), hence reality.  The support conditions follow from
\(j=(n+h)/2\) and \(k=(n-h)/2\).  Finally,
\(\e^{\ii h\pi}=(-1)^h=(-1)^n\) for every admissible \(h\), proving
\cref{eq:D-parity}.
\end{proof}

\subsection{The first three sectors}

Direct simplification of \cref{eq:Dn-definition} gives
\begin{align}
 D_1(\theta)
 &=\frac2\lambda\cos\theta,
 \label{eq:D1}\\
 D_2(\theta)
 &=-\frac1\lambda
   \left(2+\cos2\theta+2\rho\sin2\theta\right),
 \label{eq:D2}\\
 D_3(\theta)
 &=\frac1\lambda\left[
   (6-\rho^2)\cos\theta+5\rho\sin\theta
   +\left(\frac23-3\rho^2\right)\cos3\theta
   +3\rho\sin3\theta
 \right].
 \label{eq:D3}
\end{align}
Thus
\begin{align}
 X_+(y)={}&\xi+\frac{2q}{\lambda}\cos\theta
 -\frac{q^2}{\lambda}
   \left(2+\cos2\theta+2\rho\sin2\theta\right)
 \notag\\
 &+\frac{q^3}{\lambda}\left[
   (6-\rho^2)\cos\theta+5\rho\sin\theta
   +\left(\frac23-3\rho^2\right)\cos3\theta
   +3\rho\sin3\theta
 \right]
 +\Ord(q^4).
 \label{eq:first-three-sectors}
\end{align}
The leading form in the original variable is
\begin{equation}\label{eq:leading-y}
 X_+(y)=\log_{\varphi}(\sqrt5\,y)
 +\frac{\cos\!\bigl(\pi\log_{\varphi}(\sqrt5\,y)\bigr)}
 {5(\log\varphi)y^2}
 +\Ord(y^{-4}).
\end{equation}

The fourth sector is compactly represented by positive-frequency
coefficients.  If
\begin{equation}\label{eq:D-positive-frequency}
 D_n(\theta)=d_{n,0}+
 2\Re\sum_{\substack{1\le h\le n\\h\equiv n\ (2)}}
 d_{n,h}\e^{\ii h\theta},
\end{equation}
with the constant term omitted for odd \(n\), then
\begin{align}
 d_{4,0}&=-\frac9\lambda,\label{eq:d40}\\
 d_{4,2}&=-\frac{2(\lambda-\ii\omega)(2\lambda-\ii\omega)
                    (3\lambda-2\ii\omega)}{3\lambda^4},
 \label{eq:d42}\\
 d_{4,4}&=-\frac{(\lambda-4\ii\omega)(\lambda-2\ii\omega)
                    (3\lambda-4\ii\omega)}{12\lambda^4}.
 \label{eq:d44}
\end{align}
The general formula \cref{eq:dhat-nh} is usually preferable to fully
expanded trigonometric expressions at higher orders.

\begin{remark}[Amplitude and phase form]
Every conjugate pair can equivalently be written
\[
 2|\widehat D_{n,h}|q^n
 \cos\!\left(h\theta+\arg\widehat D_{n,h}\right).
\]
The phase shifts are therefore algebraic consequences of the complex
binomial coefficients; they are not independent fitting parameters.
\end{remark}

\begin{summarybox}
The outer inverse consists of the logarithmic carrier \(\xi\) plus a
convergent hierarchy of log-periodic sectors.  The sector of radial order
\(q^n=(5y^2)^{-n}\) contains precisely the harmonics
\(n,n-2,\ldots,-n\).  Formula \cref{eq:dhat-nh} gives every Fourier
coefficient directly.
\end{summarybox}
\section{Recursive construction and rigorous remainders}\label{sec:recurrence}

The closed formula is ideal for theory, but a triangular recurrence is often
more convenient for symbolic implementation.  Return to the normalized
scalar equation \cref{eq:H-equation} and seek
\begin{equation}\label{eq:delta-series-D}
 \delta(q,\theta)=\sum_{n\ge1}D_n(\theta)q^n.
\end{equation}
Suppose \(D_1,\ldots,D_{n-1}\) are known and set
\begin{equation}\label{eq:delta-less-n}
 \delta_{<n}(q,\theta)=\sum_{r=1}^{n-1}D_r(\theta)q^r.
\end{equation}

\begin{proposition}[Triangular sector recurrence]\label{prop:sector-recurrence}
For every \(n\ge1\),
\begin{equation}\label{eq:sector-recurrence}
 \boxed{
 D_n(\theta)=-\frac1\lambda[q^n]\left(
 \e^{\lambda\delta_{<n}}
 +q^2\e^{-\lambda\delta_{<n}}
 -2q\cos(\theta+\omega\delta_{<n})-1
 \right).
 }
\end{equation}
Here \([q^n]\) means extraction of the coefficient of \(q^n\), treating
\(\theta\) as an independent parameter.
\end{proposition}

\begin{proof}
Substitute
\(
 \delta=\delta_{<n}+D_nq^n+\Ord(q^{n+1})
\)
into \cref{eq:H-equation}.  In \(\e^{\lambda\delta}\), the unknown sector
contributes \(\lambda D_nq^n\).  The cosine term is multiplied by \(q\), so
its dependence on \(D_nq^n\) begins at order \(q^{n+1}\); the second
exponential is multiplied by \(q^2\), so its dependence begins still later.
Thus the coefficient equation at order \(q^n\) is linear with coefficient
\(\lambda\), yielding \cref{eq:sector-recurrence}.
\end{proof}

\subsection{Bell-polynomial form}

The recurrence can be expanded using complete exponential Bell polynomials.
For example, if
\(
 E(q)=\exp(\lambda\delta(q))=\sum_{n\ge0}E_nq^n
\), then
\begin{equation}\label{eq:Bell-exp}
 E_n=\frac1{n!}
 B_n\!\left(1!\lambda D_1,2!\lambda D_2,\
 \ldots,n!\lambda D_n\right),
\end{equation}
where \(B_n\) is the complete Bell polynomial.  The Fourier cone is
preserved under all products appearing in these Bell polynomials.  A similar
expansion of
\(
 \exp(\ii\omega\delta)
\)
handles the cosine through
\[
 \cos(\theta+\omega\delta)
 =\Re\!\left(\e^{\ii\theta}\e^{\ii\omega\delta}\right).
\]
This gives a purely algebraic, finite computation at each order.

The closed coefficient formula and the triangular recurrence are useful for
different purposes.  Formula \cref{eq:Cjk} exposes arithmetic and asymptotic
structure; recurrence \cref{eq:sector-recurrence} is easy to truncate,
simplify, and evaluate numerically without manipulating generalized
binomial functions of large complex arguments.

\subsection{Residual certification}

For \(N\ge0\), let
\begin{equation}\label{eq:delta-N}
 \delta_N(q,\theta)=\sum_{n=1}^ND_n(\theta)q^n,
 \qquad \delta_0=0.
\end{equation}
By construction,
\begin{equation}\label{eq:residual-order}
 H(\delta_N;q,\theta)=\Ord(q^{N+1}).
\end{equation}
The derivative of \(H\) is
\begin{equation}\label{eq:H-delta}
 H_\delta(\delta;q,\theta)=
 \lambda\e^{\lambda\delta}
 -\lambda q^2\e^{-\lambda\delta}
 +2\omega q\sin(\theta+\omega\delta).
\end{equation}
For \(|q|\) and \(|\delta|\) small, this is bounded away from zero
uniformly in real \(\theta\).

\begin{theorem}[Uniform remainder estimate]\label{thm:uniform-remainder}
For each \(N\ge0\), there exist \(q_N>0\) and \(C_N>0\) such that, for
real \(\theta\) and \(0<q<q_N\),
\begin{equation}\label{eq:uniform-remainder}
 \left|
 \delta(q,\theta)-\delta_N(q,\theta)
 \right|\le C_Nq^{N+1}.
\end{equation}
Equivalently,
\begin{equation}\label{eq:uniform-remainder-y}
 X_+(y)-\left(\xi+\sum_{n=1}^ND_n(\theta)(5y^2)^{-n}\right)
 =\Ord(y^{-2N-2}).
\end{equation}
For every fixed ordinary derivative order \(m\ge0\),
\begin{equation}\label{eq:derivative-remainder}
 \frac{\dd^m}{\dd y^m}
 \left[
 X_+(y)-\xi-\sum_{n=1}^ND_n(\theta)(5y^2)^{-n}
 \right]
 =\Ord(y^{-2N-2-m}).
\end{equation}
\end{theorem}

\begin{proof}
Equation \cref{eq:residual-order} and the uniform lower bound for
\(|H_\delta|\) imply \cref{eq:uniform-remainder} by the mean-value theorem
on the real branch, or by the analytic inverse-function theorem with a
uniform compact parameter \(\theta\in\R/2\pi\mathbb Z\).  Since
\(q=(5y^2)^{-1}\), this gives \cref{eq:uniform-remainder-y}.  The exact
solution is analytic in \((X,Y)\) near the origin, so the normalized
remainder is analytic in \(q\) with smooth periodic dependence on
\(\theta\).  Applying
\[
 y\frac{\dd}{\dd y}=-2q\frac{\partial}{\partial q}
 +2\rho\frac{\partial}{\partial\theta}
\]
preserves radial order.  Converting Euler derivatives to ordinary
derivatives adds one factor \(y^{-1}\) per differentiation, proving
\cref{eq:derivative-remainder}.
\end{proof}

\subsection{Newton acceleration in the transseries algebra}

A second computational route is Newton iteration applied to
\cref{eq:H-equation}:
\begin{equation}\label{eq:Newton-delta}
 \delta^{\mathrm{new}}
 =\delta-\frac{H(\delta;q,\theta)}{H_\delta(\delta;q,\theta)}.
\end{equation}
All operations are performed in the truncated algebra \(\Talg/(q^{N+1})\).
If the current approximation has residual valuation \(r\), the next Newton
iterate ordinarily has residual valuation at least \(2r\).  Starting from
\(\delta=0\), repeated doubling therefore computes \(N\) sectors in
\(\Ord(\log N)\) Newton stages, assuming fast truncated multiplication.
The triangular recurrence remains simpler for hand calculation, while
Newton iteration is attractive for very high order.

\begin{algorithm}[Certified high-order inverse evaluation]\label{alg:evaluate}
For a prescribed \(y\), sector order \(N\), and target precision:
\begin{enumerate}
\item compute \(R=5y^2\), \(q=R^{-1}\),
      \(\xi=\log R/\lambda\), and \(\theta=\pi\xi\);
\item generate \(D_1,\ldots,D_N\) either from
      \cref{eq:dhat-nh} or \cref{eq:sector-recurrence};
\item evaluate \(x_N=\xi+\sum_{n=1}^ND_n(\theta)q^n\);
\item compute the exact residual
      \(r_N=\Fcal(x_N)-y\), or preferably the squared residual in
      \cref{eq:R-equation-x};
\item convert the residual to an a posteriori error estimate using a lower
      bound for \(\Fcal'\) near \(x_N\), or perform one full-precision
      Newton correction.
\end{enumerate}
\end{algorithm}

\section{Coefficient structure}\label{sec:coefficient-structure}

\subsection{Diagonal, edge, and mean sectors}

Formula \cref{eq:dhat-nh} distinguishes several natural families.
The edge harmonics \(h=\pm n\) correspond to \((j,k)=(n,0)\) and
\((0,n)\).  The central harmonic \(h=0\) exists only for even \(n=2m\)
and corresponds to \(j=k=m\).  Since \(s_{m,m}=2m\), its coefficient is
particularly simple:
\begin{equation}\label{eq:central-coeff}
 \widehat D_{2m,0}
 =-\frac1{2m\lambda}\binom{2m}{m}^2.
\end{equation}
The first values are
\[
 -\frac2\lambda,
 \quad -\frac9\lambda,
 \quad -\frac{200}{3\lambda},
 \quad -\frac{1225}{2\lambda},
 \quad\ldots.
\]
These constants are the angular means of the even sectors.

\subsection{Hypergeometric and elliptic resummation of the mean}

For a periodic function, define
\begin{equation}\label{eq:angular-mean}
 \mean{f}=\frac1{2\pi}\int_0^{2\pi}f(\theta)\dd\theta.
\end{equation}
Averaging \cref{eq:all-orders-real} and using
\cref{eq:central-coeff} gives
\begin{equation}\label{eq:mean-series}
 \mean{\delta}(q)
 =-\frac1\lambda\sum_{m\ge1}
 \frac1{2m}\binom{2m}{m}^2q^{2m}.
\end{equation}
This subseries has the closed hypergeometric form
\begin{equation}\label{eq:mean-4F3}
 \boxed{
 \mean{\delta}(q)=
 -\frac{2q^2}{\lambda}\,
 {}_4F_3\!\left(
 \begin{matrix}1,1,\tfrac32,\tfrac32\\2,2,2\end{matrix};16q^2
 \right).
 }
\end{equation}
Let \(\mathbf K(k)\) denote the complete elliptic integral in modulus
notation,
\begin{equation}\label{eq:K-modulus}
 \mathbf K(k)=\int_0^{\pi/2}
 \frac{\dd t}{\sqrt{1-k^2\sin^2t}}.
\end{equation}
The classical generating function
\begin{equation}\label{eq:central-square-K}
 \sum_{m\ge0}\binom{2m}{m}^2q^{2m}
 =\frac2\pi\mathbf K(4q)
\end{equation}
then implies
\begin{equation}\label{eq:mean-derivative-K}
 q\frac{\dd}{\dd q}\mean{\delta}(q)
 =-\frac1\lambda\left(\frac2\pi\mathbf K(4q)-1\right).
\end{equation}
The mean sector has radius \(|q|<1/4\).  The complete transseries has a
smaller uniform radius because off-diagonal harmonics grow more rapidly.

\begin{remark}
Averaging over the logarithmic phase is not the same operation as evaluating
at a typical large \(y\): the physical angle \(\theta=\rho\log R\) keeps
rotating.  Nevertheless, the mean isolates a distinguished exactly
summable subsector and reveals a direct elliptic-integral connection.
\end{remark}

\subsection{Large-order growth and the exact uniform radius}

Set \(j=np\), \(k=n(1-p)\) heuristically and define
\begin{equation}\label{eq:sigma-p}
 \sigma(p)=1-\ii\rho(2p-1),
 \qquad 0\le p\le1.
\end{equation}
With principal logarithms and the convention \(0\log0=0\), put
\begin{align}
 g_\rho(p)={}&\Re\!\left[
 2\sigma\Logp\sigma
 -(\sigma-p)\Logp(\sigma-p)
 -(\sigma-1+p)\Logp(\sigma-1+p)
 \right]
 \notag\\
 &-p\log p-(1-p)\log(1-p).
 \label{eq:g-rho}
\end{align}
The function is continuous and symmetric under \(p\mapsto1-p\).

\begin{theorem}[Coefficient growth and uniform convergence radius]\label{thm:uniform-radius}
Let
\begin{equation}\label{eq:Gstar-qstar}
 G_\rho=\exp\!\left(\max_{0\le p\le1}g_\rho(p)\right),
 \qquad q_*=G_\rho^{-1}.
\end{equation}
Then
\begin{equation}\label{eq:C-growth}
 \limsup_{n\to\infty}
 \max_{0\le j\le n}|C_{j,n-j}|^{1/n}=G_\rho,
\end{equation}
and, with the supremum norm on continuous \(2\pi\)-periodic functions,
\begin{equation}\label{eq:D-growth}
 \limsup_{n\to\infty}\|D_n\|_\infty^{1/n}=G_\rho.
\end{equation}
Consequently the Banach-valued power series
\(\sum_{n\ge1}D_n(\theta)q^n\) has exact radius \(q_*\) for uniform
convergence in \(\theta\).
\end{theorem}

\begin{proof}
Write the generalized binomial coefficients in \cref{eq:Cjk} as Gamma
quotients and apply Stirling's formula with \(j/n\to p\).  The terms
proportional to \(n\log n\) and \(n\) cancel, leaving precisely
\(g_\rho(p)\) as the exponential rate.  Uniform Stirling estimates on
closed subintervals, followed by continuous endpoint limits, give
\cref{eq:C-growth}.

Every \(C_{j,n-j}\) is a Fourier coefficient of \(D_n\), hence
\[
 \max_j|C_{j,n-j}|\le\|D_n\|_\infty.
\]
Conversely,
\[
 \|D_n\|_\infty
 \le\sum_{j=0}^n|C_{j,n-j}|
 \le(n+1)\max_j|C_{j,n-j}|.
\]
Taking \(n\)-th roots proves \cref{eq:D-growth}; the Cauchy--Hadamard
formula gives the radius.
\end{proof}

For the Fibonacci value \(\rho=\pi/(2\log\varphi)\), numerical maximization
gives two symmetric maximizers,
\begin{equation}\label{eq:pstar-numeric}
 p_*=0.0934868615594596\ldots,
 \qquad 1-p_*=0.9065131384405404\ldots,
\end{equation}
and
\begin{align}
 G_\rho&=10.1243628948610529\ldots,
 \label{eq:G-numeric}\\
 q_*&=0.09877164720237184\ldots.
 \label{eq:qstar-numeric}
\end{align}
Therefore the Fourier-sector expansion converges absolutely and uniformly
whenever
\begin{equation}\label{eq:y-uniform-threshold}
 q=(5y^2)^{-1}<q_*;
 \qquad\text{in particular, whenever}\qquad
 y>1.422980175185941\ldots.
\end{equation}
This sufficient-and-sharp threshold concerns uniform convergence as a
power series with values in \(C(\R/2\pi\mathbb Z)\).  Special phases may
have larger pointwise domains; no such enlargement is needed for the
large-argument theory.

The dominant large-order coefficients are not the angular means.  The mean
has per-radial-order growth \(4^n\), while the full norm grows like
\(G_\rho^n\) with \(G_\rho\approx10.1244\).  The maximizing normalized
frequency is
\[
 \frac{|h|}{n}=|2p_*-1|
 =0.813026276881081\ldots,
\]
so high-order growth is driven by off-center harmonics.

\begin{summarybox}
The inverse transseries is a genuinely convergent expansion, not merely a
formal asymptotic device.  Its exact uniform radial radius is determined by
a one-dimensional entropy-like maximization derived from complex Stirling
asymptotics.  For every integer Fibonacci checkpoint beginning with
\(y=F_3=2\), the physical value \(q=(5y^2)^{-1}\) lies comfortably inside
the uniform disk.
\end{summarybox}
\section{Numerical verification}\label{sec:numerics}

The explicit series can be tested without any ambiguity by substituting its
partial sums into the original function.  Define
\begin{equation}\label{eq:XN-def}
 X_{+,N}(y)=\xi+\sum_{n=1}^ND_n(\theta)q^n,
 \qquad X_{+,0}(y)=\xi.
\end{equation}
The following values were computed with high-precision arithmetic.  The
entries \(E_N\) are absolute errors
\(
 |X_{+,N}(y)-X_+(y)|
\).  At \(y=2\) and \(y=5\), the exact inverse values are integers because
\(2=F_3\) and \(5=F_5\).

\begin{table}[H]
\centering
\caption{Exact inverse values and absolute errors of sector truncations.}
\label{tab:numerical-errors}
\small
\begin{tabular}{@{}r l@{}}
\toprule
\(y\) & exact \(X_+(y)\)\\
\midrule
2   & \(3\)\\
5   & \(5\)\\
10  & \(6.45779309028821164945\ldots\)\\
100 & \(11.24218977013280389416\ldots\)\\
\bottomrule
\end{tabular}

\medskip
\begin{tabular}{@{}r r r r r@{}}
\toprule
\(y\) & \(E_0\) & \(E_1\) & \(E_2\) & \(E_3\)\\
\midrule
2
& \(1.1269603\!\times10^{-1}\)
& \(1.5236037\!\times10^{-2}\)
& \(2.9618737\!\times10^{-3}\)
& \(2.2607545\!\times10^{-3}\)\\
5
& \(1.6827815\!\times10^{-2}\)
& \(2.2634536\!\times10^{-4}\)
& \(1.8596489\!\times10^{-5}\)
& \(9.6297844\!\times10^{-7}\)\\
10
& \(5.4518532\!\times10^{-4}\)
& \(1.1352239\!\times10^{-5}\)
& \(1.5425214\!\times10^{-7}\)
& \(2.1471593\!\times10^{-9}\)\\
100
& \(3.0101615\!\times10^{-5}\)
& \(3.5619587\!\times10^{-9}\)
& \(2.9715886\!\times10^{-13}\)
& \(5.7169682\!\times10^{-18}\)\\
\bottomrule
\end{tabular}

\medskip
\begin{tabular}{@{}r r r r@{}}
\toprule
\(y\) & \(E_4\) & \(E_6\) & \(E_8\)\\
\midrule
2
& \(3.7691613\!\times10^{-4}\)
& \(1.2868198\!\times10^{-4}\)
& \(1.5141006\!\times10^{-5}\)\\
5
& \(2.9886773\!\times10^{-8}\)
& \(1.8233102\!\times10^{-11}\)
& \(6.5804309\!\times10^{-13}\)\\
10
& \(3.0170906\!\times10^{-11}\)
& \(6.0001812\!\times10^{-15}\)
& \(1.1659689\!\times10^{-18}\)\\
100
& \(4.3013834\!\times10^{-21}\)
& \(1.4868585\!\times10^{-28}\)
& \(2.1225057\!\times10^{-36}\)\\
\bottomrule
\end{tabular}
\end{table}

The convergence is phase dependent, so errors need not decrease
monotonically at every adjacent order.  For example, at \(y=2\) the
order-10 and order-11 errors are approximately
\(1.1640\times10^{-6}\) and \(1.4004\times10^{-6}\), respectively.
The all-orders estimate concerns the magnitude of
the remainder after a fixed truncation as \(y\to\infty\), while convergence
for a fixed admissible \(y\) follows from the power-series theorem.  By
\(y=10\), each additional sector already contributes roughly two more
decimal digits; at \(y=100\), the leading carrier is accurate to five
decimal places and a few sectors give extremely high precision.

\subsection{Integer checkpoints}

For every integer \(n\ge2\), \cref{prop:integer-interpolation} and
\cref{prop:outer-monotone} imply
\begin{equation}\label{eq:integer-checkpoint}
 X_+(F_n)=n.
\end{equation}
For \(n\ge3\), one has \(q=(5F_n^2)^{-1}\le1/20<q_*\), so the
uniform convergence theorem applies.  Substituting \(y=F_n\) into
\cref{eq:all-orders-complex} therefore yields the convergent identity:
\begin{equation}\label{eq:integer-transseries-identity}
 n=\log_{\varphi}(\sqrt5\,F_n)
 +\sum_{\substack{j,k\ge0\\j+k\ge1}}
 C_{jk}(5F_n^2)^{-s_{jk}}.
\end{equation}
This is not a tautological restatement of Binet's formula at the level of
individual sectors.  It decomposes the exact difference between \(n\) and
the logarithmic Binet estimate into a convergent hierarchy of log-periodic
complex-power corrections.

At integer \(n\), Binet's formula gives
\[
 \sqrt5\,F_n=\varphi^n\left(1-(-1)^n\varphi^{-2n}\right),
\]
so
\begin{equation}\label{eq:baseline-integer-error}
 \log_{\varphi}(\sqrt5\,F_n)
 =n+\frac1{\log\varphi}
 \log\!\left(1-(-1)^n\varphi^{-2n}\right).
\end{equation}
The leading sector in \cref{eq:leading-y} cancels the first term of the
logarithm in \cref{eq:baseline-integer-error}; all higher sectors complete
the exact cancellation.  This provides a stringent symbolic and numerical
check on signs and phase conventions.

\section{Wolfram Language implementation}\label{sec:wl}

The user-supplied interpolation is already expressed in Wolfram Language,
so it is natural to give an implementation in the same notation.  The
following code evaluates the exact coefficient formula and groups it by
radial order.

\begin{lstlisting}[caption={Closed coefficients and truncated inverse transseries.},label={lst:wl-coeff}]
ClearAll[lam, rho, s, c, dn, inverseSeries];

lam = 2 ArcCsch[2];                 (* = 2 Log[GoldenRatio] *)
rho = Pi/lam;

s[j_Integer?NonNegative, k_Integer?NonNegative] :=
  (1 - I rho) j + (1 + I rho) k;

c[j_Integer?NonNegative, k_Integer?NonNegative] /; j + k > 0 :=
  (-1)^(j + k + 1)/(lam s[j, k]) *
    Binomial[s[j, k], j] Binomial[s[j, k], k];

dn[n_Integer?Positive, theta_] :=
  FullSimplify[
    ComplexExpand[
      Sum[c[j, n - j] Exp[I (2 j - n) theta], {j, 0, n}],
      TargetFunctions -> {Re, Im}
    ],
    Assumptions -> Element[theta, Reals]
  ];

inverseSeries[y_, nmax_Integer?NonNegative] := Module[
  {r = 5 y^2, q, xi, theta},
  q = 1/r;
  xi = Log[r]/lam;
  theta = Pi xi;
  xi + Sum[dn[n, theta] q^n, {n, 1, nmax}]
];
\end{lstlisting}

For numerical validation, one may compare against a high-precision root of
the original function.

\begin{lstlisting}[caption={Direct outer inverse and residual check.},label={lst:wl-root}]
ClearAll[f, outerInverse, check];

f[x_] := Sqrt[2/5] Sqrt[Cosh[lam x] - Cos[Pi x]];

outerInverse[y_?NumericQ, wp_: 80] := Module[{x0},
  x0 = N[Log[5 y^2]/lam, wp];
  x /. FindRoot[
    f[x] == SetPrecision[y, wp],
    {x, x0},
    WorkingPrecision -> wp,
    AccuracyGoal -> Floor[wp/2],
    PrecisionGoal -> Floor[wp/2]
  ]
];

check[y_?NumericQ, nmax_Integer?NonNegative, wp_: 80] := Module[
  {xa, xt},
  xa = N[inverseSeries[SetPrecision[y, wp], nmax], wp];
  xt = outerInverse[y, wp];
  <|
    "Approximation" -> xa,
    "DirectRoot" -> xt,
    "AbsoluteError" -> Abs[xa - xt],
    "FunctionResidual" -> Abs[f[xa] - y]
  |>
];
\end{lstlisting}

The recurrence \cref{eq:sector-recurrence} can also be implemented by
repeatedly forming a truncated \texttt{Series}.  The closed formula is less
stateful and naturally parallelizes over pairs \((j,k)\); the recurrence is
usually faster after simplification when one needs all sectors up to a large
common order.

\subsection{Stable numerical evaluation of high-order coefficients}

For moderate \(n\), Wolfram Language's exact generalized
\texttt{Binomial} expression is reliable.  At large \(n\), direct products
can overflow or suffer cancellation in machine arithmetic.  Three remedies
are standard:
\begin{enumerate}
\item evaluate with arbitrary precision and use \texttt{LogGamma} to form
      Gamma quotients in logarithmic scale;
\item pair \((j,k)\) with \((k,j)\) before summing, so every contribution is
      explicitly real;
\item use the triangular recurrence or Newton iteration in the real Fourier
      basis, avoiding large intermediate complex binomials.
\end{enumerate}
For a requested numerical value rather than symbolic coefficients, a few
transseries sectors followed by one Newton step is typically the most
stable and economical procedure.
\part{General inversion theory}

\section{Finite exponential products}\label{sec:general-products}

The Fibonacci factorization is a special case of a broad class whose
inverses admit equally explicit coefficients.  Let \(m\ge1\), let
\(\lambda\ne0\), and consider a chosen sector in which
\(\Re(\lambda x)\to+\infty\).  Fix parameters
\begin{equation}\label{eq:general-params}
 c_r,\beta_r\in\C,
 \qquad \alpha_r\in\C,
 \qquad \Re\alpha_r>0,
 \qquad r=1,\ldots,m.
\end{equation}
Branches of powers are chosen by analytic continuation from factors near
\(1\).

Consider the normalized product map
\begin{equation}\label{eq:general-product-map}
 R=\e^{\lambda x}
 \prod_{r=1}^m
 \left(1-c_r\e^{-\lambda\alpha_r x}\right)^{\beta_r}.
\end{equation}
A nonzero constant prefactor \(A\) causes no difficulty: replace \(R\) by
\(R/A\).

\subsection{Fixed-point normalization}

Set
\begin{equation}\label{eq:general-uq}
 u=\e^{-\lambda x},
 \qquad q=R^{-1}.
\end{equation}
Then \cref{eq:general-product-map} is
\begin{equation}\label{eq:general-u-equation}
 Ru=\prod_{r=1}^m(1-c_ru^{\alpha_r})^{\beta_r}.
\end{equation}
Introduce
\begin{equation}\label{eq:general-AQV}
 A_r=c_ru^{\alpha_r},
 \qquad
 Q_r=c_rq^{\alpha_r},
 \qquad
 V=\frac uq=\prod_{r=1}^m(1-A_r)^{\beta_r}.
\end{equation}
Then
\begin{equation}\label{eq:general-fixed-system}
 A_r=Q_rV^{\alpha_r},
 \qquad r=1,\ldots,m,
\end{equation}
and the inverse correction is
\begin{equation}\label{eq:general-delta}
 x=\frac{\log R}{\lambda}-\frac1\lambda\log V.
\end{equation}
This is the same architecture as the Fibonacci problem, with an arbitrary
number of factors and arbitrary exponents.

For a multi-index \(\multi n=(n_1,\ldots,n_m)\in\Nzero^m\), write
\begin{equation}\label{eq:S-multi}
 |\multi n|=\sum_{r=1}^mn_r,
 \qquad
 S_{\multi n}=\sum_{r=1}^m\alpha_rn_r,
 \qquad
 \multi c^{\multi n}=\prod_{r=1}^mc_r^{n_r}.
\end{equation}
Because every \(\Re\alpha_r>0\), one has
\(\Re S_{\multi n}>0\) for every nonzero \(\multi n\).

\begin{theorem}[Closed inverse of a finite exponential product]\label{thm:general-product}
For \(|R|\) sufficiently large in the selected sector, the inverse branch
of \cref{eq:general-product-map} satisfying
\(x\sim\lambda^{-1}\log R\) is
\begin{equation}\label{eq:general-product-inverse}
 \boxed{
 x(R)=\frac{\log R}{\lambda}
 +\sum_{\multi n\in\Nzero^m\setminus\{\multi0\}}
 \frac{(-1)^{|\multi n|+1}}{\lambda S_{\multi n}}
 \left[
   \prod_{r=1}^m
   \binom{\beta_rS_{\multi n}}{n_r}c_r^{n_r}
 \right]
 R^{-S_{\multi n}}.
 }
\end{equation}
The multi-series converges near \(R=\infty\) in its natural variables
\(Q_r=c_rR^{-\alpha_r}\), and it is a complete asymptotic transseries when
ordered by \(\Re S_{\multi n}\).
\end{theorem}

\begin{proof}
Apply \cref{thm:good} to \cref{eq:general-fixed-system}, where
\[
 \Phi_r(\multi A)=V(\multi A)^{\alpha_r},
 \qquad
 V(\multi A)=\prod_{s=1}^m(1-A_s)^{\beta_s}.
\]
The Good matrix is
\begin{equation}\label{eq:general-Good-matrix}
 M_{rs}=\delta_{rs}
 +\alpha_r\frac{\beta_sA_s}{1-A_s}.
\end{equation}
It is \(I+\multi a\multi b^{\mathsf T}\), where
\(a_r=\alpha_r\) and
\(b_s=\beta_sA_s/(1-A_s)\).  The matrix determinant lemma gives
\begin{equation}\label{eq:general-Delta}
 \Delta(\multi A)=
 1+\sum_{r=1}^m\frac{\alpha_r\beta_rA_r}{1-A_r}.
\end{equation}

Define
\begin{equation}\label{eq:Pn-general}
 P_{\multi n}(t)=
 [\multi A^{\multi n}]V(\multi A)^t\Delta(\multi A).
\end{equation}
The constant part of the determinant contributes
\begin{equation}\label{eq:Bn-general}
 B_{\multi n}(t)=
 (-1)^{|\multi n|}
 \prod_{r=1}^m\binom{\beta_rt}{n_r}.
\end{equation}
For the summand indexed by \(r\) in \cref{eq:general-Delta}, use
\[
 \beta_r\binom{\beta_rt-1}{n_r-1}
 =\frac{n_r}{t}\binom{\beta_rt}{n_r}.
\]
Its sign is opposite to \cref{eq:Bn-general}, so summing all determinant
terms gives
\begin{equation}\label{eq:Pn-collapse}
 P_{\multi n}(t)=B_{\multi n}(t)
 \left(1-\frac{S_{\multi n}}t\right).
\end{equation}

For an auxiliary exponent \(\tau\), Good inversion gives
\begin{equation}\label{eq:general-Vtau}
 [\multi Q^{\multi n}]V(\multi A(\multi Q))^\tau
 =P_{\multi n}(\tau+S_{\multi n}).
\end{equation}
Differentiation at \(\tau=0\) yields the coefficient of \(\log V\).  The
factor in parentheses in \cref{eq:Pn-collapse} vanishes at
\(t=S_{\multi n}\), hence
\begin{equation}\label{eq:general-logV-coeff}
 [\multi Q^{\multi n}]\log V
 =\frac{(-1)^{|\multi n|}}{S_{\multi n}}
 \prod_{r=1}^m\binom{\beta_rS_{\multi n}}{n_r}.
\end{equation}
Now use \cref{eq:general-delta} and
\(
 \multi Q^{\multi n}=\multi c^{\multi n}R^{-S_{\multi n}}
\) to obtain \cref{eq:general-product-inverse}.  Analytic convergence near
\(\multi Q=\multi0\) follows from the implicit-function theorem.
\end{proof}

\subsection{Checks and special cases}

\begin{example}[One factor]\label{ex:one-factor}
For
\[
 R=\e^{\lambda x}(1-c\e^{-\lambda\alpha x})^\beta,
\]
\cref{eq:general-product-inverse} reduces to
\begin{equation}\label{eq:one-factor-inverse}
 x=\frac{\log R}{\lambda}
 +\sum_{n\ge1}
 \frac{(-1)^{n+1}}{\lambda\alpha n}
 \binom{\beta\alpha n}{n}c^nR^{-\alpha n}.
\end{equation}
When \(\alpha=\beta=1\), the original equation is
\(R=\e^{\lambda x}-c\), whose exact inverse is
\[
 x=\frac1\lambda\log(R+c)
 =\frac{\log R}{\lambda}
 +\frac1\lambda\sum_{n\ge1}\frac{(-1)^{n+1}}n
 \left(\frac cR\right)^n.
\]
Formula \cref{eq:one-factor-inverse} reproduces it because
\(\binom nn=1\).
\end{example}

\begin{example}[Fibonacci conjugate pair]\label{ex:fib-product}
Take \(m=2\),
\[
 c_1=c_2=1,
 \qquad \beta_1=\beta_2=1,
 \qquad \alpha_1=1-\ii\rho,
 \qquad \alpha_2=1+\ii\rho.
\]
Then \(S_{(j,k)}=s_{jk}\), and
\cref{eq:general-product-inverse} becomes exactly
\cref{eq:all-orders-complex}.
\end{example}

\begin{example}[Purely real corrections]\label{ex:real-product}
If all \(\alpha_r,c_r,\beta_r\) are real and
\(\alpha_r>0\), every exponent \(S_{\multi n}\) is real.  The result is an
ordinary generalized power series in real powers of \(R^{-1}\), with no
logarithmic oscillation.  Log-periodicity appears precisely when nonreal
exponents occur.
\end{example}

\subsection{Reality, periodicity, and quasiperiodicity}

Suppose the parameter list is closed under complex conjugation, with
conjugate values of \(c_r\), \(\alpha_r\), and \(\beta_r\).  Multi-indices
are then paired by conjugation, and the inverse transseries is real on the
positive real axis.  A monomial with
\(
 S_{\multi n}=a-\ii b
\) takes the form
\begin{equation}\label{eq:general-log-periodic-monomial}
 R^{-S_{\multi n}}=R^{-a}\e^{\ii b\log R}.
\end{equation}
If all frequencies \(b\) lie in an integer multiple of one basic frequency,
the coefficients are periodic functions of \(\log R\).  If several basic
frequencies are rationally independent, the coefficient functions are
quasiperiodic on a finite-dimensional torus.  The same inversion theorem
handles both cases without modification; only the presentation of the
monomial group changes.

\subsection{Support and local finiteness}

Let
\begin{equation}\label{eq:Gamma-support}
 \Gamma=\left\{
 S_{\multi n}:\multi n\in\Nzero^m
 \right\}.
\end{equation}
Because \(\min_r\Re\alpha_r>0\), for every \(A>0\) there are only finitely
many multi-indices satisfying \(\Re S_{\multi n}\le A\).  Thus the support
is locally finite with respect to real asymptotic order.  Multiplication,
analytic substitution, and coefficient extraction are therefore finite at
every prescribed order.  This is the elementary finite-grid version of the
well-based support condition used in general transseries and Hahn-series
calculus.

\section{Arbitrary analytic perturbations}\label{sec:analytic-perturbations}

The product form is special because its coefficients collapse to a product
of binomial factors.  The normalization itself needs only analyticity.
Let
\begin{equation}\label{eq:Psi-map}
 R=\e^{\lambda x}
 \Psi\!\left(
 c_1\e^{-\lambda\alpha_1x},\ldots,
 c_m\e^{-\lambda\alpha_mx}
 \right),
\end{equation}
where \(\Psi\) is analytic near the origin and
\begin{equation}\label{eq:Psi-zero}
 \Psi(\multi0)=1,
 \qquad \Re\alpha_r>0.
\end{equation}
With \(u=\e^{-\lambda x}\), \(q=R^{-1}\),
\[
 A_r=c_ru^{\alpha_r},
 \qquad Q_r=c_rq^{\alpha_r},
 \qquad V=\frac uq,
\]
the inverse equation becomes
\begin{equation}\label{eq:Psi-fixed}
 V=\Psi(\multi A),
 \qquad
 A_r=Q_rV^{\alpha_r},
 \qquad
 x=\frac{\log R}{\lambda}-\frac1\lambda\log V.
\end{equation}

\begin{theorem}[Analytic perturbation inversion]\label{thm:Psi-inversion}
Under \cref{eq:Psi-zero}, there is a unique analytic germ
\(V(\multi Q)=1+\Ord(\multi Q)\) solving \cref{eq:Psi-fixed}.  Hence the
inverse branch has a convergent transseries
\begin{equation}\label{eq:Psi-inverse-series}
 x(R)=\frac{\log R}{\lambda}
 +\sum_{\multi n\ne\multi0}a_{\multi n}
 \multi c^{\multi n}R^{-S_{\multi n}},
 \qquad S_{\multi n}=\sum_r\alpha_rn_r.
\end{equation}
Its coefficients admit the extraction formula
\begin{equation}\label{eq:Psi-coeff-formula}
 \boxed{
 a_{\multi n}= -\frac1\lambda
 \left.\frac{\dd}{\dd t}\right|_{t=S_{\multi n}}
 [\multi A^{\multi n}]
 \Psi(\multi A)^t
 \det\!\left(
  \delta_{rs}-\alpha_rA_s\frac{\partial}{\partial A_s}
                    \log\Psi(\multi A)
 \right).
 }
\end{equation}
\end{theorem}

\begin{proof}
The equations
\(
 A_r=Q_r\Psi(\multi A)^{\alpha_r}
\)
have the solution \(\multi A=\multi0\) at \(\multi Q=\multi0\), and their
Jacobian with respect to \(\multi A\) is the identity there.  The analytic
implicit-function theorem gives the unique germ.  Apply
\cref{thm:good} with
\(
 \Phi_r=\Psi^{\alpha_r}
\) and auxiliary function \(\Psi^\tau\).  The coefficient is the expression
inside \cref{eq:Psi-coeff-formula} evaluated at
\(t=\tau+S_{\multi n}\); differentiation at \(\tau=0\), followed by the
factor \(-1/\lambda\), proves the result.
\end{proof}

The product theorem is recovered by setting
\(
 \Psi(\multi A)=\prod_r(1-A_r)^{\beta_r}
\).  Its special rank-one Jacobian and binomial factorization are what turn
the general coefficient extraction \cref{eq:Psi-coeff-formula} into the
closed product \cref{eq:general-product-inverse}.

\subsection{A recursive alternative}

For general \(\Psi\), a total-degree recursion may be simpler than
\cref{eq:Psi-coeff-formula}.  Write
\[
 V=1+W,
 \qquad W\in(\multi Q),
\]
and iterate
\begin{equation}\label{eq:Psi-picard}
 W=\Psi\!\left(Q_1(1+W)^{\alpha_1},\ldots,
 Q_m(1+W)^{\alpha_m}\right)-1
\end{equation}
by total degree.  The right-hand side at degree \(N\) depends only linearly
on the unknown degree-\(N\) part of \(W\), and the coefficient of that part
is zero at the origin, so the recursion is triangular.  Newton iteration in
the complete multi-adic algebra gives quadratic order-doubling when many
coefficients are required.

\subsection{Uniform analytic estimates}

Choose a polydisc \(|Q_r|\le r_r\) on which the implicit solution is
analytic and remains away from zeros and branch cuts of \(\Psi\).  Cauchy's
estimate gives
\begin{equation}\label{eq:multi-Cauchy}
 |a_{\multi n}|\le M\prod_{r=1}^mr_r^{-n_r}
\end{equation}
for a suitable \(M\).  Along a physical ray,
\(|Q_r|=|c_r||R|^{-\Re\alpha_r}\) times a bounded angular factor.  This
immediately supplies uniform truncation estimates by any half-plane of the
support,
\(
 \Re S_{\multi n}\le A
\), and justifies termwise differentiation on smaller polydiscs.

\begin{summarybox}
A dominant exponential plus any analytic combination of finitely many
subdominant exponential monomials has a convergent inverse transseries on
the additive semigroup generated by their exponents.  Product perturbations
are the exactly solvable subclass in which multivariate Good inversion
collapses to generalized binomial coefficients.
\end{summarybox}
\section{Mixed carriers and nested transseries}\label{sec:mixed-carriers}

The preceding theory assumes that the dominant carrier is a pure
exponential.  Many special functions instead have a carrier of the form
\begin{equation}\label{eq:mixed-carrier}
 C(x)=A\e^{\lambda x}x^\mu(\log x)^\nu
 \left(1+\sum_{r\ge1}\frac{P_r(\log x)}{x^r}\right),
\end{equation}
possibly multiplied by exponentially small factors.  Its inverse already
belongs to a polynomial--logarithmic or Lambert-type transseries.  The
log-periodic sectors developed above can be placed over that slower
coefficient field.

\subsection{Carrier-first normalization}

Consider
\begin{equation}\label{eq:carrier-Psi}
 Y=C(x)\,
 \Psi\!\left(c_1\e^{-\zeta_1x},\ldots,
             c_m\e^{-\zeta_mx}\right),
 \qquad \Psi(\multi0)=1,
 \qquad \Re\zeta_r>0.
\end{equation}
Let \(x_0=x_0(Y)\) be a selected inverse branch of the carrier,
\begin{equation}\label{eq:x0-carrier}
 C(x_0)=Y,
\end{equation}
and write \(x=x_0+\delta\).  Define
\begin{equation}\label{eq:carrier-Q}
 Q_r=c_r\e^{-\zeta_rx_0}.
\end{equation}
The exact correction equation is
\begin{equation}\label{eq:carrier-correction-exact}
 \log C(x_0+\delta)-\log C(x_0)
 +\log\Psi\!\left(Q_1\e^{-\zeta_1\delta},\ldots,
                  Q_m\e^{-\zeta_m\delta}\right)=0.
\end{equation}
If
\begin{equation}\label{eq:kappa-def}
 \kappa(x_0)=\frac{\dd}{\dd x}\log C(x)\bigg|_{x=x_0}
\end{equation}
is nonzero, the implicit-function theorem gives a unique correction
\(\delta=\Ord(\multi Q)\).  Its first term is
\begin{equation}\label{eq:carrier-first-correction}
 \delta=-\frac{\log\Psi(\multi Q)}{\kappa(x_0)}
 +\Ord(\|\multi Q\|^2).
\end{equation}
For a product perturbation
\(
 \Psi=\prod_r(1-Q_r)^{\beta_r}
\), this becomes
\begin{equation}\label{eq:carrier-first-product}
 \delta=\frac1{\kappa(x_0)}
 \sum_{r=1}^m\beta_rQ_r+\Ord(\|\multi Q\|^2).
\end{equation}

The higher coefficients involve derivatives of \(\log C\) at \(x_0\) and
therefore belong to the differential coefficient field generated by
\(x_0(Y)\).  The small monomials are
\(
 \e^{-\zeta_rx_0(Y)}
\), which may themselves be powers of \(Y\) multiplied by powers of
\(\log Y\) and iterated logarithms.  This is the natural nested-transseries
construction: invert the slow carrier first, then perform a finite-grid
exponential reversion over its coefficient field.

\subsection{The carrier \texorpdfstring{\(\e^{\lambda x}x^\mu\)}{exp(lambda x) times x to the mu} and Lambert \texorpdfstring{\(W\)}{W}}

Take
\begin{equation}\label{eq:exp-power-carrier}
 C(x)=A\e^{\lambda x}x^\mu.
\end{equation}
For \(\mu\ne0\), a carrier inverse is
\begin{equation}\label{eq:x0-W}
 x_0(Y)=\frac\mu\lambda
 W_k\!\left(
   \frac\lambda\mu\left(\frac YA\right)^{1/\mu}
 \right),
\end{equation}
with the branch \(W_k\) selected to match the desired sector.  Indeed,
raising \(\e^{\lambda x}x^\mu=Y/A\) to the power \(1/\mu\) reduces the
equation to the defining form \(we^w=z\).  When \(\mu=0\),
\begin{equation}\label{eq:x0-pure-exp}
 x_0(Y)=\frac1\lambda\log\frac YA.
\end{equation}
The logarithmic derivative is
\begin{equation}\label{eq:kappa-exp-power}
 \kappa(x_0)=\lambda+\frac\mu{x_0}.
\end{equation}
Consequently the first product correction is
\begin{equation}\label{eq:first-W-sector}
 \delta=\frac{
   \sum_r\beta_rc_r\e^{-\zeta_rx_0}}
 {\lambda+\mu/x_0}
 +\Ord\!\left(\max_r|\e^{-\zeta_rx_0}|^2\right).
\end{equation}
Expanding \(W_k\) at infinity turns the coefficients into
polynomials in \(\log Y\) and \(\log\log Y\), while the factors
\(\e^{-\zeta_rx_0}\) become complex powers of \(Y\) decorated by slow
logarithmic powers.  Thus polynomial--logarithmic and log-periodic
transseries are not competing descriptions: the former supplies the slow
coefficient field, and the latter supplies the exponentially separated
sectors.

\subsection{A formal implicit-function principle over a slow field}

Let \(\mathscr K\) be a differential field of slow transseries in a large
variable \(Y\), closed under the operations needed to describe \(x_0(Y)\)
and the derivatives of \(\log C\).  Let \(\mathfrak m\) be the ideal
generated by the small monomials \(Q_r=\e^{-\zeta_rx_0}\).  If
\(\kappa(x_0)\) is a unit of \(\mathscr K\), then
\cref{eq:carrier-correction-exact} has a unique solution
\begin{equation}\label{eq:delta-slow-field}
 \delta\in\mathfrak m\,
 \mathscr K[[Q_1,\ldots,Q_m]].
\end{equation}
The proof is the same recursive implicit-function argument as before: the
linearized coefficient of \(\delta\) is the unit \(\kappa(x_0)\), and every
nonlinear operation raises total \(Q\)-degree.  This principle is the
appropriate abstraction for inverses of Gamma-, Barnes-, hyperfactorial-,
and Lambert-type carriers perturbed by subdominant exponential products.

\begin{warningbox}
One should not expand the slow carrier and the exponentially small sectors
in a single undifferentiated list of terms.  The correct organization is
lexicographic: first by exponential sector, then by the internal
polynomial--logarithmic scale of its coefficient.  Mixing the two orderings
can destroy local finiteness and makes remainder statements ambiguous.
\end{warningbox}

\section{Complex inverse branches}\label{sec:complex-branches}

The real interpolation uses a square root and modulus interpretation, but
the squared equation
\begin{equation}\label{eq:complex-squared}
 R=\e^{\lambda x}+\e^{-\lambda x}-2\cos(\omega x)
\end{equation}
is entire in \(x\) and can be inverted locally in the complex plane.  Fix a
branch of \(\Logp R\).  The dominant equation \(\e^{\lambda x}=R\) has
logarithmic branches
\begin{equation}\label{eq:x0m}
 x_{0,m}=\frac{\Logp R+2\pi\ii m}{\lambda},
 \qquad m\in\mathbb Z.
\end{equation}
Define
\begin{equation}\label{eq:ZmXYm}
 Z_m=\e^{\ii\omega x_{0,m}},
 \qquad
 X_m=R^{-1}Z_m,
 \qquad
 Y_m=R^{-1}Z_m^{-1}.
\end{equation}
Writing \(x=x_{0,m}+\delta\), the normalized equation factors as
\begin{align}
 1={}&\e^{\lambda\delta}
 \left(1-X_m\e^{(-\lambda+\ii\omega)\delta}\right)
 \left(1-Y_m\e^{(-\lambda-\ii\omega)\delta}\right).
 \label{eq:complex-factor-delta}
\end{align}
With \(V=\e^{-\lambda\delta}\), this is again
\[
 A=X_mV^\alpha,
 \qquad B=Y_mV^{\bar\alpha},
 \qquad V=(1-A)(1-B).
\]

\begin{proposition}[Right-exponential complex branches]\label{prop:complex-branches}
Whenever \(|X_m|\) and \(|Y_m|\) are sufficiently small, there is a unique
solution near \(x_{0,m}\), and it is
\begin{equation}\label{eq:complex-branch-series}
 x_m(R)=x_{0,m}+
 \sum_{\substack{j,k\ge0\\j+k\ge1}}C_{jk}X_m^jY_m^k,
\end{equation}
with the same universal coefficients \(C_{jk}\) as in
\cref{eq:Cjk}.
\end{proposition}

The condition involves both \(|X_m|\) and \(|Y_m|\).  Since
\(
 |Z_m|=\exp(-2\pi\rho m)
\) for positive real \(R\) and the principal logarithm, moving far in one
direction of the branch index makes one generator small but the conjugate
generator large.  Thus \cref{eq:complex-branch-series} describes precisely
the complex branches for which the \(\e^{\lambda x}\) balance remains
dominant.  Other complex solutions may be governed by a different dominant
balance and require a separate Newton polygon or transseries chart.

The real positive branch is \(m=0\), with
\(Z_0=\e^{\ii\theta}\), and the real negative branch comes from the exact
even symmetry \(x\mapsto-x\), not from changing \(m\) in the right-exponential
family.

\section{Interpretation of the log-periodic sectors}\label{sec:interpretation}

\subsection{Discrete scale covariance}

The phase is
\begin{equation}\label{eq:phase-logR}
 \theta=\rho\log R=\pi\log_{\varphi}(\sqrt5\,y).
\end{equation}
A shift \(\theta\mapsto\theta+2\pi\) corresponds to
\begin{equation}\label{eq:scale-period}
 R\mapsto R\e^{2\pi/\rho}=R\e^{2\lambda}=R\varphi^4,
 \qquad
 y\mapsto\varphi^2y.
\end{equation}
Thus every coefficient function \(D_n(\theta)\) repeats under the discrete
scale factor \(\varphi^2\) in \(y\), while its radial amplitude acquires the
expected factor \(\varphi^{-4n}\).  This is covariance rather than exact
periodicity of the whole correction: each sector transforms with its own
weight.

Because \(D_n(\theta+\pi)=(-1)^nD_n(\theta)\), the half-period scale
\(y\mapsto\varphi y\) changes the sign of odd sectors and preserves even
ones, while multiplying the \(n\)-th radial amplitude by \(\varphi^{-2n}\).
This parity mirrors the alternating term in Binet's formula.

\subsection{Complex dimensions}

The exponents
\begin{equation}\label{eq:complex-dimensions}
 s_{jk}=n-\ii\rho h
\end{equation}
form a two-dimensional lattice projected onto the complex plane.  Their
real parts determine decay, and their imaginary parts determine logarithmic
frequency.  In Mellin-transform language, such exponents play the role of
complex poles or complex dimensions.  Here they arise algebraically before
any transform is taken: they are generated by the two conjugate factors in
\cref{eq:factorized-u}.  The inverse operation preserves the additive
semigroup of these exponents but renormalizes their amplitudes through
complex generalized binomial coefficients.

\subsection{No Stokes ambiguity on the positive real branch}

The expansion is convergent in a nontrivial domain, so its coefficients do
not hide an exponentially small ambiguity.  On the positive real branch,
there is no choice of transseries parameter analogous to a Stokes constant:
analytic continuation from \(q=0\) fixes the solution uniquely.  Stokes
phenomena may enter only after one enlarges the problem to competing complex
dominant balances or to divergent slow-carrier expansions.  This distinction
is useful: the log-periodic sectors are small relative to inverse logarithms,
but they are not nonperturbative ambiguities.  They are ordinary convergent
Taylor sectors in the lifted variables \((X,Y)\).

\section{Further deductions and research directions}\label{sec:future}

The exact formula opens several directions that are not visible from a few
low-order coefficients.

\subsection{Pointwise singularity curves}

The exact uniform radius \(q_*\) is controlled by coefficient growth in the
supremum norm over all phases.  For a fixed physical phase trajectory
\(\theta=-\rho\log q\), cancellation among harmonics may postpone the first
singularity.  Determining the maximal analytic continuation of
\(q\mapsto\delta(q,-\rho\log q)\) requires solving simultaneously
\begin{equation}\label{eq:singularity-system}
 H(\delta;q,\theta)=0,
 \qquad
 H_\delta(\delta;q,\theta)=0,
 \qquad
 \theta=-\rho\log q
\end{equation}
with compatible logarithm branches.  The resulting singularities should
govern exponentially accurate large-order asymptotics of the physical
one-variable coefficient combinations.

\subsection{Asymptotics near the dominant coefficient rays}

The maximizers \(p_*,1-p_*\) in \cref{eq:pstar-numeric} identify dominant
rays \(j/n\to p_*\) in coefficient space.  A two-dimensional saddle-point
analysis of \(C_{j,n-j}\), followed by a local Fourier summation in \(j\),
should produce an asymptotic approximation to \(D_n(\theta)\) uniform in
\(\theta\).  Depending on phase, the two conjugate saddles may reinforce or
cancel.  This would refine the root test into a complete estimate of the
prefactor, phase, and transition regions.

\subsection{Arithmetic at Fibonacci arguments}

Identity \cref{eq:integer-transseries-identity} evaluates the convergent
complex-power series at \(R=5F_n^2\) to an integer.  The individual terms
are generally transcendental-looking, involving
\(
 \exp(\ii\rho h\log(5F_n^2))
\), yet their full sum has exact arithmetic meaning.  One may seek finite
transformations, functional equations, or resummations that make the Binet
alternation manifest term by term.  The half-period parity
\cref{eq:D-parity} suggests separating even and odd radial sectors before
specializing to Fibonacci arguments.

\subsection{Other Lucas sequences}

For a Lucas sequence with characteristic roots \(r\) and \(s\), a modulus
interpolation of
\(
 (r^x-s^x)/(r-s)
\)
usually leads, after squaring, to a dominant exponential plus conjugate or
real subdominant exponentials.  When the squared expression factorizes into
a finite product in \(u=\e^{-\lambda x}\),
\cref{thm:general-product} applies directly.  The nature of the inverse
sectors is controlled by the ratios of logarithms of the characteristic
roots: real ratios yield generalized powers, while imaginary parts yield
log-periodic phases.

\subsection{Formalization prospects}

The essential proof can be decomposed into formalizable algebraic pieces:
\begin{enumerate}
\item the identity \(\arccsch2=\log\varphi\) and the squared Binet formula;
\item monotonicity of the outer branch on \([2,\infty)\);
\item the factorization \cref{eq:factorized-u};
\item a formal multivariate power-series implicit-function theorem;
\item Good's coefficient formula and the rank-one determinant calculation;
\item finite coefficient identities in generalized binomial polynomials;
\item analytic convergence and residual bounds on an explicit neighborhood.
\end{enumerate}
The algebraic coefficient theorem is particularly suitable for
formalization because every coefficient at a fixed multi-index reduces to a
finite polynomial identity; Gamma functions are unnecessary if generalized
binomial coefficients are defined by falling factorials.

\subsection{Conjectural sharp physical radius}

It is natural to distinguish three radii: the bivariate domain of the lifted
germ, the uniform Fourier radius \(q_*\), and the radius seen along the
physical logarithmic spiral.  Numerical analytic continuation could test
the following working conjecture.

\begin{quote}
\textbf{Conjecture.}
The first singularity encountered when the positive physical
logarithmic-spiral branch is continued inward from arbitrarily small
positive \(q\) is a square-root branch point produced by the system
\cref{eq:singularity-system}; its modulus is at least \(q_*\), with strict
inequality unless a maximizing Fourier phase aligns with the logarithmic
spiral.
\end{quote}

The conjecture is deliberately stated as a program rather than a theorem.
The uniform radius proved in \cref{thm:uniform-radius} is exact in its stated
Banach-space sense and does not depend on this pointwise question.
\section{Conclusion}

The real-argument Fibonacci modulus has an inverse problem whose apparent
mixture of an exponential, a reciprocal exponential, and a trigonometric
term conceals a simple product structure.  After squaring, setting
\(u=\e^{-\lambda x}\), and normalizing by \(R=5y^2\), the entire problem is
encoded by
\[
 Ru=(1-u^{1-\ii\rho})(1-u^{1+\ii\rho}).
\]
The two conjugate factors lift the inverse to a bivariate analytic germ.
Multivariate Lagrange--Good inversion then yields the exact coefficient
formula \cref{eq:Cjk}; restriction to the physical logarithmic spiral gives
the real Fourier sectors \cref{eq:dhat-nh}.  The result is simultaneously a
complete asymptotic transseries, a convergent expansion, and an effective
high-precision algorithm.

Several features are worth retaining as a general template.  One should
first identify the dominant carrier, then rewrite subdominant terms as a
finite set of small exponential monomials, and finally lift those monomials
to independent variables before applying formal or analytic inversion.
Conjugate complex exponents become periodic functions of the logarithm;
rationally independent imaginary parts become quasiperiodic coefficient
functions.  Product perturbations have closed generalized-binomial
coefficients, while arbitrary analytic perturbations retain an explicit
Good-determinant extraction formula.  When the carrier itself requires a
Lambert or polynomial--logarithmic inverse, the same exponential-sector
calculus operates over the resulting slow coefficient field.

For the specific interpolation, the two real large-value inverse branches
are therefore completely described:
\[
 X_-(y)=-X_+(y),
 \qquad
 X_+(y)=\log_{\varphi}(\sqrt5\,y)
 +\sum_{n\ge1}D_n\!\left(
   \pi\log_{\varphi}(\sqrt5\,y)
 \right)(5y^2)^{-n}.
\]
Every coefficient is explicit, every fixed truncation has a certified
remainder order, and the full Fourier-sector series has a rigorously
identified uniform convergence radius.

\appendix

\section{Factorized Fourier coefficients through order six}\label{app:coefficients}

For convenience, define positive-frequency coefficients by
\begin{equation}\label{eq:appendix-positive}
 D_n(\theta)=d_{n,0}+
 2\Re\sum_{\substack{1\le h\le n\\h\equiv n\ (2)}}
 d_{n,h}\e^{\ii h\theta},
\end{equation}
where \(d_{n,0}\) is present only for even \(n\).  In all formulas below,
\(\omega=\pi\).  The universal expression is
\begin{equation}\label{eq:appendix-dnh-general}
 d_{n,h}=\frac{(-1)^{n+1}}{\lambda(n-\ii\rho h)}
 \binom{n-\ii\rho h}{(n+h)/2}
 \binom{n-\ii\rho h}{(n-h)/2}.
\end{equation}
The following factorizations avoid unnecessary expanded polynomials.

\begingroup
\small
\begin{align*}
 d_{1,1}&=\frac1\lambda;\\[0.35em]
 d_{2,0}&=-\frac2\lambda,
 &d_{2,2}&=-\frac{\lambda-2\ii\omega}{2\lambda^2};\\[0.35em]
 d_{3,1}&=\frac{(2\lambda-\ii\omega)(3\lambda-\ii\omega)}
                    {2\lambda^3},
 &d_{3,3}&=\frac{(\lambda-3\ii\omega)(2\lambda-3\ii\omega)}
                    {6\lambda^3};\\[0.35em]
 d_{4,0}&=-\frac9\lambda,\\
 d_{4,2}&=-\frac{2(\lambda-\ii\omega)(2\lambda-\ii\omega)
                    (3\lambda-2\ii\omega)}{3\lambda^4},\\
 d_{4,4}&=-\frac{(\lambda-4\ii\omega)(\lambda-2\ii\omega)
                    (3\lambda-4\ii\omega)}{12\lambda^4};\\[0.35em]
 d_{5,1}&=\frac{(3\lambda-\ii\omega)(4\lambda-\ii\omega)^2
                    (5\lambda-\ii\omega)}{12\lambda^5},\\
 d_{5,3}&=\frac{(\lambda-\ii\omega)(2\lambda-3\ii\omega)
                    (4\lambda-3\ii\omega)(5\lambda-3\ii\omega)}
                    {8\lambda^5},\\
 d_{5,5}&=\frac{(\lambda-5\ii\omega)(2\lambda-5\ii\omega)
                    (3\lambda-5\ii\omega)(4\lambda-5\ii\omega)}
                    {120\lambda^5};\\[0.35em]
 d_{6,0}&=-\frac{200}{3\lambda},\\
 d_{6,2}&=-\frac{(2\lambda-\ii\omega)(3\lambda-2\ii\omega)
                    (3\lambda-\ii\omega)(5\lambda-2\ii\omega)^2}
                    {12\lambda^6},\\
 d_{6,4}&=-\frac{2(\lambda-2\ii\omega)(\lambda-\ii\omega)
                    (3\lambda-4\ii\omega)(3\lambda-2\ii\omega)
                    (5\lambda-4\ii\omega)}{15\lambda^6},\\
 d_{6,6}&=-\frac{(\lambda-6\ii\omega)(\lambda-3\ii\omega)
                    (\lambda-2\ii\omega)(2\lambda-3\ii\omega)
                    (5\lambda-6\ii\omega)}{60\lambda^6}.
\end{align*}
\endgroup

The negative-frequency coefficients are
\(d_{n,-h}=\overline{d_{n,h}}\).  These expressions also provide convenient
unit tests for a symbolic implementation of \cref{eq:dhat-nh}.

\section{Direct derivation of the first sectors}\label{app:direct-sectors}

This appendix verifies \cref{eq:D1,eq:D2,eq:D3} without multivariate
inversion.  Substitute
\[
 \delta=D_1q+D_2q^2+D_3q^3+\Ord(q^4)
\]
into \cref{eq:H-equation}.  The coefficient of \(q\) is
\begin{equation}\label{eq:direct-q1}
 \lambda D_1-2\cos\theta=0,
\end{equation}
so \(D_1=2\cos\theta/\lambda\).

At order \(q^2\), one obtains
\begin{equation}\label{eq:direct-q2}
 \lambda D_2+\frac{\lambda^2}{2}D_1^2+1
 +2\omega D_1\sin\theta=0.
\end{equation}
Substituting \(D_1\) and using elementary double-angle identities yields
\[
 D_2=-\frac1\lambda
 \left(2+\cos2\theta+2\rho\sin2\theta\right).
\]

At order \(q^3\), the coefficient equation is
\begin{equation}\label{eq:direct-q3}
 \lambda D_3
 +\lambda^2D_1D_2
 +\frac{\lambda^3}{6}D_1^3
 -\lambda D_1
 +2\omega D_2\sin\theta
 +\omega^2D_1^2\cos\theta=0.
\end{equation}
After substituting the first two sectors and reducing products of sines and
cosines to harmonics \(1\) and \(3\), this becomes \cref{eq:D3}.  The
calculation illustrates both the triangular nature of the scalar recurrence
and the rapid growth of algebraic complexity that the closed coefficient
formula avoids.

\section{Residue proof of the Lagrange--Good formula}\label{app:good-proof}

For completeness, we outline a formal-residue proof of
\cref{thm:good}.  Write
\begin{equation}\label{eq:Q-change}
 Q_r=\frac{A_r}{\Phi_r(\multi A)}.
\end{equation}
The hypothesis \(\Phi_r(\multi0)\ne0\) makes this a formally invertible
change of variables.  Coefficient extraction is a multivariate residue:
\begin{equation}\label{eq:coeff-residue}
 [\multi Q^{\multi n}]f(\multi A(\multi Q))
 =\Res_{\multi Q=\multi0}
 f(\multi A(\multi Q))
 \prod_{r=1}^mQ_r^{-n_r-1}\dd Q_r.
\end{equation}
The Jacobian of \cref{eq:Q-change} is
\begin{equation}\label{eq:Q-Jacobian}
 \det\!\left(\frac{\partial Q_r}{\partial A_s}\right)
 =\left(\prod_{r=1}^m\Phi_r^{-1}\right)
 \det\!\left(
  \delta_{rs}-A_r\frac{\partial}{\partial A_s}
                         \log\Phi_r
 \right).
\end{equation}
By the identity \(\det(I-UV)=\det(I-VU)\), the determinant in
\cref{eq:Q-Jacobian} equals
\begin{equation}\label{eq:Good-det-equivalent}
 \det\!\left(
  \delta_{rs}-A_s\frac{\partial}{\partial A_s}
                         \log\Phi_r
 \right)=\Delta(\multi A).
\end{equation}
Meanwhile,
\begin{equation}\label{eq:Q-powers-change}
 \prod_rQ_r^{-n_r-1}
 =\prod_rA_r^{-n_r-1}\Phi_r^{n_r+1}.
\end{equation}
Multiplying \cref{eq:Q-Jacobian} and
\cref{eq:Q-powers-change} cancels one factor \(\Phi_r\) for each \(r\),
leaving
\[
 f(\multi A)\prod_r\Phi_r^{n_r}\Delta(\multi A)
 \prod_rA_r^{-n_r-1}\dd A_r.
\]
Its residue is exactly the coefficient on the right-hand side of
\cref{eq:Good-formula}.  The same proof is valid analytically on sufficiently
small tori or formally in the ring of multivariate Laurent series.

\section{Notation crosswalk}\label{app:notation}

\begin{longtable}{@{}p{0.21\textwidth}p{0.72\textwidth}@{}}
\toprule
Symbol & Meaning\\
\midrule
\endhead
\(\varphi\) & Golden ratio \((1+\sqrt5)/2\).\\
\(\lambda\) & Dominant exponential rate \(2\log\varphi\).\\
\(\omega\) & Oscillatory frequency \(\pi\).\\
\(\rho\) & Frequency ratio \(\omega/\lambda\).\\
\(R\) & Squared output variable \(5y^2\).\\
\(q\) & Small radial variable \(R^{-1}\).\\
\(\xi\) & Leading inverse carrier \(\lambda^{-1}\log R
=\log_{\varphi}(\sqrt5\,y)\).\\
\(\theta\) & Logarithmic phase \(\rho\log R=\pi\xi\).\\
\(u\) & Exponential inverse variable \(\e^{-\lambda x}\).\\
\(\alpha,\bar\alpha\) & Conjugate exponents \(1\mp\ii\rho\).\\
\(s_{jk}\) & Complex radial exponent \(j+k-\ii\rho(j-k)\).\\
\(C_{jk}\) & Bivariate coefficient in \cref{eq:Cjk}.\\
\(D_n(\theta)\) & Real Fourier polynomial multiplying \(q^n\).\\
\(S_{\multi n}\) & General multi-factor exponent
\(\sum_r\alpha_rn_r\).\\
\bottomrule
\end{longtable}

\begin{thebibliography}{99}

\bibitem{Corless1996}
R.~M. Corless, G.~H. Gonnet, D.~E.~G. Hare, D.~J. Jeffrey, and D.~E. Knuth,
``On the Lambert \(W\) function,''
\emph{Advances in Computational Mathematics} \textbf{5} (1996), 329--359.

\bibitem{FlajoletGourdonDumas1995}
P.~Flajolet, X.~Gourdon, and P.~Dumas,
``Mellin transforms and asymptotics: harmonic sums,''
\emph{Theoretical Computer Science} \textbf{144} (1995), 3--58.

\bibitem{Gessel1987}
I.~M. Gessel,
``A combinatorial proof of the multivariable Lagrange inversion formula,''
\emph{Journal of Combinatorial Theory, Series A} \textbf{45} (1987),
178--195.

\bibitem{Gessel2016}
I.~M. Gessel,
``Lagrange inversion,''
\emph{Journal of Combinatorial Theory, Series A} \textbf{144} (2016),
212--249; arXiv:1609.05988.

\bibitem{Good1960}
I.~J. Good,
``Generalizations to several variables of Lagrange's expansion, with
applications to stochastic processes,''
\emph{Proceedings of the Cambridge Philosophical Society} \textbf{56}
(1960), 367--380.

\bibitem{GKP1994}
R.~L. Graham, D.~E. Knuth, and O.~Patashnik,
\emph{Concrete Mathematics}, 2nd ed., Addison--Wesley, 1994.

\bibitem{Koshy2001}
T.~Koshy,
\emph{Fibonacci and Lucas Numbers with Applications}, Wiley, 2001.

\bibitem{Olver1997}
F.~W.~J. Olver,
\emph{Asymptotics and Special Functions}, A K Peters, 1997 reprint.

\bibitem{ProveItTransseries}
V.~Reshetnikov,
``Series and transseries,'' in the \emph{ProveIt} Fabius-function research
drafts, 2026,
\url{https://github.com/VladimirReshetnikov/ProveIt/tree/main/Analysis/FabiusFunction/docs/semi-formalized-research-frontiers/drafts/series-and-transseries}.

\bibitem{Sornette1998}
D.~Sornette,
``Discrete-scale invariance and complex dimensions,''
\emph{Physics Reports} \textbf{297} (1998), 239--270.

\bibitem{vanderHoeven2006}
J.~van der Hoeven,
\emph{Transseries and Real Differential Algebra},
Lecture Notes in Mathematics 1888, Springer, 2006.

\end{thebibliography}

\end{document}
